\documentclass[a4paper]{article}
\usepackage{import}
\input{../../../preamble.sty}

\begin{document}

\begin{titlepage}
	\begin{center}
		\vspace*{1cm}

		\Huge
		\textbf{Basi di dati}

		\vspace{0.5cm}
		\LARGE
		UniVR - Dipartimento di Informatica

		\vspace{1.5cm}

		\textbf{Fabio Irimie}

		\vfill


		\vspace{0.8cm}


		2° Semestre 2025/2026

	\end{center}
\end{titlepage}


\tableofcontents
\pagebreak

\section{Tecnologia dei DBMS}
Un DBMS (Database Management System) è un software in grado di gestire collezioni di dati
che siano \textbf{grandi, condivise e persistenti} assicurando \textbf{affidabilità e
privatezza}.
\begin{itemize}
  \item Grandi: i dati devono essere organizzati in modo efficiente per poter essere gestiti
    in grandi quantità
  \item Condivise: i dati devono essere accessibili da più utenti contemporaneamente
  \item Persistenti: i dati devono essere memorizzati in modo permanente, anche dopo la
    chiusura del programma
  \item Affidabilità: i dati devono essere protetti da errori e guasti del sistema
  \item Privatezza: i dati devono essere protetti da accessi non autorizzati
\end{itemize}

Un DBMS deve essere anche \textbf{efficiente} ed \textbf{efficace}:
\begin{itemize}
  \item Efficiente: il DBMS deve essere in grado di gestire grandi quantità di dati inbasitex
    modo rapido e con un utilizzo minimo delle risorse
  \item Efficace: il DBMS deve essere in grado di soddisfare le esigenze degli utenti
    in modo completo e accurato
\end{itemize}
Il DBMS si appoggia sul filesystem estendendo le sue funzionalità.

\subsection{Architettura transazionale di un DBMS}
Un DBMS è composto da tre livelli:
\begin{figure}[H]
  \centering
  \includegraphics[width=0.8\textwidth]{architettura_3_livelli}
  \caption{Architettura a 3 livelli di un DBMS}
\end{figure}
\noindent
Un DBMS basato sul \textbf{modello relazionale} è nella maggior parte dei casi anche
un \textbf{sistema transazionale}, cioè fornisce un meccanismo per la definizione ed
esecuzione di transazioni:
\begin{figure}[H]
  \centering
  \includegraphics[width=0.8\textwidth]{architettura_transazionale}
  \caption{Architettura transazionale di un DBMS}
\end{figure}

\subsection{Transazione}
\begin{definition}
  È \textbf{un’unità} di lavoro svolta da un programma applicativo (che interagisce con una base
  di dati) per la quale si vogliono garantire proprietà di correttezza, robustezza e
  isolamento.
\end{definition}
La principale caratteristica di una transazione è che deve essere \textbf{atomica}, cioè
deve essere eseguita completamente o non essere eseguita affatto. \textbf{Non sono
ammesse esecuzioni parziali}.

\subsubsection{Sintassi}
La sintassi per definire una transazione è la seguente:
\begin{lstlisting}
<transazione> -> begin transaction
                 <programma>
                 end transaction

<programma> -> < <istruzione> | commit work | rollback work >
               { < <istruzione> | commit work | rollback work > }
\end{lstlisting}
\begin{itemize}
  \item La transazione va a buon fine all'esecuzione di un \texttt{commit work}.
  \item La transazione non ha alcun effetto se viene eseguito un \texttt{rollback work}.
\end{itemize}

\vspace{1em}
\noindent
Una transazione è ben formata se:
\begin{itemize}
  \item Inizia con un \texttt{begin transaction}.
  \item Termina con un \texttt{end transaction}.
  \item La sua esecuzione comporta il raggiungimento di un \texttt{commit work} o di un
    \texttt{rollback work} e dopo il \texttt{commit/rollback} non si eseguono altri
    accessi alla base di dati.
\end{itemize}

\begin{example}
  Un esempio di transazione ben formata è il seguente:
\begin{lstlisting}
begin transaction;
  update CONTO set saldo = saldo - 1200
    where filiale = '005' and numero = 15;
  update CONTO set saldo = saldo + 1200
    where filiale = '005' and numero = 205;
  commit work;
end transaction;
\end{lstlisting}
\end{example}

\subsubsection{Proprietà delle transazionei}
Una transazione ha 4 proprietà fondamentali, note come ACID:
\begin{itemize}
  \item \textbf{Atomicità} (\textbf{A}tomicity):
    Una transazione è una unità di esecuzione \textbf{indivisibile}. O viene eseguita
    completamente o non viene eseguita affatto.

    \vspace{1em}
    \noindent
    Implicazioni:
    \begin{itemize}
      \item Se una transazione viene interrotta prima del commit, il lavoro fin qui
        eseguito dalla transazione deve essere disfatto ripristinando la situazione in
        cui si trovava la base di dati prima dell’inizio della transazione.

      \item Se una transazione viene interrotta dopo l’esecuzione del commit (commit
        eseguito con successo), il sistema deve assicurare che la transazione abbia
        effetto sulla base di dati.
    \end{itemize}

  \item \textbf{Consistenza} (\textbf{C}onsistency):
    L’esecuzione di una transazione non deve violare i vincoli di integrità.

    \vspace{1em}
    \noindent
    Implicazioni:
    \begin{itemize}
      \item Verifica immediata:
        \begin{itemize}
          \item Fatta nel corso della transazione
          \item Viene abortita solo l’ultima operazione e il sistema restituisce
            all’applicazione una segnalazione d’errore
          \item L’applicazione può quindi reagire alla violazione.
        \end{itemize}

      \item Verifica differita:
        \begin{itemize}
          \item Al commit se un vincolo di integrità viene violato la transazione viene
            abortita senza possibilità da parte dell’applicazione di reagire alla
            violazione.
        \end{itemize}
    \end{itemize}

  \item \textbf{Isolamento} (\textbf{I}solation):
    L’esecuzione di una transazione deve essere indipendente dalla contemporanea
    esecuzione di altre transazioni.

    \vspace{1em}
    \noindent
    Implicazioni:
    \begin{itemize}
      \item Il rollback di una transazione non deve creare rollback a catena di altre
        transazioni che si trovano in esecuzione contemporaneamente.

      \item Il sistema deve regolare l’esecuzione concorrente con meccanismi di controllo
        dell’accesso alle risorse.
    \end{itemize}

  \item \textbf{Persistenza} (\textbf{D}urability):
    L’effetto di una transazione che ha eseguito il commit non deve andare perso.

    \vspace{1em}
    \noindent
    Implicazioni:
    \begin{itemize}
      \item Il sistema deve essere in grado, in caso di guasto, di garantire gli effetti
        delle transazioni che al momento del guasto avevano già eseguito un commit.
    \end{itemize}
\end{itemize}
Un DBMS che gestisce transazioni dovrebbe garantire per ogni transazione che esegue
tutte queste proprietà.

\subsection{Architettura di di un DBMS}
L'architettura di un DBMS è composta da \textbf{moduli} che si occupano di gestire
le diverse funzionalità del sistema.

\subsubsection{Vista ad alto livello}
\begin{figure}[H]
  \centering
  \includegraphics[width=0.8\textwidth]{architettura_alto_livello}
  \caption{Vista ad alto livello dell'architettura di un DBMS}
\end{figure}

\subsubsection{Modulo di gestione delle transazioni}
Questo modulo si occupa di trasformare una transazione in un insieme di operazioni
sulla base di dati, garantendo le proprietà ACID.
\begin{figure}[H]
  \centering
  \includegraphics[width=0.8\textwidth]{modulo_gestione_transazioni}
  \caption{Modulo di gestione delle transazioni di un DBMS}
\end{figure}

I singoli moduli che compongono questo modulo sono:
\begin{itemize}
  \item \textbf{Gestore delle interrogazioni}:
    Riceve in ingresso una query SQL
    \begin{figure}[H]
      \centering
      \includegraphics[width=0.8\textwidth]{gestore_interrogazioni}
      \caption{Gestore delle interrogazioni di un DBMS}
    \end{figure}

  \item \textbf{Gestore dei metodi di accesso}: Determina quali sono le 
    pagine della memoria secondaria di cui si ha bisogno, ma capisce anche se
    è possibile accederci.
    \begin{figure}[H]
      \centering
      \includegraphics[width=0.8\textwidth]{gestore_metodi_accesso}
      \caption{Gestore dei metodi di accesso di un DBMS}
    \end{figure}

  \item \textbf{Gestore dell'affidabilità}:
    Carica in memoria principale le pagine della memoria
    secondaria di cui si ha bisogno, e le mantiene in memoria finché non sono più
    necessarie.
    \begin{figure}[H]
      \centering
      \includegraphics[width=0.8\textwidth]{gestore_affidabilita}
      \caption{Gestore del buffer di un DBMS}
    \end{figure}

  \item \textbf{Gestore del buffer}
  \item \textbf{Gestore della memoria secondaria}
  \item \textbf{Gestore della concorrenza}
\end{itemize}

\vspace{1em}
\noindent
I moduli che contribuiscono a garantire le proprietà ACID sono:
\begin{itemize}
  \item \textbf{Gestore dei metodi di accesso}: garantisce la consistenza
  \item \textbf{Gestore dell'esecuzione concorrente}: garantisce l'isolamento e l'atomicità
  \item \textbf{Gestore dell'affidabilità}: garantisce l'atomicità e la persistenza
\end{itemize}

\section{Linguaggio di interrogazione SQL}
SQL (Structured Query Language) è il linguaggio di riferimento per le basi di dati
relazionali. SQL è un linguaggio dichiarativo composto da 3 tipi di istruzioni:
\begin{itemize}
\item \textbf{Data Definition Language} (DDL): linguaggio per la definizione della
  struttura della base di dati (es. \texttt{CREATE}, \texttt{ALTER}, \texttt{DROP})

\item \textbf{Data Manipulation Language} (DML): linguaggio per manipolare i dati
  come inserimento, aggiornamento e cancellazione (es. \texttt{INSERT}, \texttt{UPDATE},
  \texttt{DELETE})

\item \textbf{Data Query Language} (DQL): linguaggio per interrogare i dati (es.
  \texttt{SELECT})
\end{itemize}
Esistono varie versioni di SQL, in questo corso ci baseremo su SQL-92 (SQL-2), che è una versione
standard del linguaggio SQL.

\vspace{1em}
\noindent
La parte di SQL dedicata alle interrogazioni fa parte dei DQL (Data Query Language). SQL esprime le
interrogazioni in modo dichiarativo (specifica l’obiettivo e non la procedura per ottenerlo,
come l’algebra relazionale). Una parte del DBMS, detta query optimizer, traduce di
fatto l’interrogazione SQL nel linguaggio procedurale interno del sistema, che è nascosto
all’utente. SQL permette quindi di descrivere le interrogazioni in modo astratto ed ad
alto livello.

\subsection{PostgreSQL}
È un Relational Database Management System (RDBMS) con alcune funzionalità orientate
agli oggetti e verrà utilizzato all'interno del corso. Il modello di interazione con
le basi di dati è \textbf{Client-Server}.

\subsection{Fondamenti del DDL}
\subsubsection{Identificatori}
Un identificatore è una sequenza di caratteri che identifica un elemento del database
come ad esempio una tabella, un attributo, un vincolo, ecc. Gli identificatori in SQL
devono rispettare le seguenti regole:
\begin{itemize}
  \item Devono iniziare con una lettera (UTF-8) o un underscore (\texttt{\_})
  \item Possono contenere lettere, numeri (0-9) e underscore (\texttt{\_})
  \item Non possono essere parole chiave riservate di SQL
\end{itemize}
Gli identificatori sono \textbf{case-insensitive}, quindi \texttt{idPersona} e
\texttt{idpersona} sono considerati lo stesso identificatore.

\subsubsection{Domini elementari}
I domini elementari sono i tipi predefiniti di SQL:
\begin{itemize}
  \item \textbf{Tipo carattere}:
    Rappresenta stringhe di lunghezza variabile secondo la seguente sintassi:
\begin{lstlisting}[language=SQL, numbers=none]
CHARACTER [VARYING][(lunghezza)]
\end{lstlisting}
    \begin{itemize}
      \item \texttt{CHARACTER}: carattere singolo

      \item \texttt{CHARACTER(20)}: stringa di lunghezza fissa di 20 caratteri.
        Se la stringa è più corta di 20 caratteri, viene riempita con spazi a destra.

      \item \texttt{CHARACTER VARYING(20)}: stringa di lunghezza variabile con un massimo
        di 20 caratteri. Se la stringa è più corta di 20 caratteri, viene memorizzata
        senza spazi aggiuntivi.

      \item \texttt{TEXT} (estensione PostgreSQL): stringa di lunghezza variabile senza
        un limite.
    \end{itemize}
    Le stringhe sono delimitate da apici singoli (\texttt{'}).

  \item \textbf{Tipo booleano}:
    Rappresenta valori booleani (vero o falso) più lo stato nullo secondo la seguente
    sintassi:
\begin{lstlisting}[language=SQL, numbers=none]
BOOLEAN
\end{lstlisting}
    La rappresentazione dei valori booleani può essere:
    \begin{itemize}
      \item \texttt{TRUE} o \texttt{FALSE}
      \item \texttt{t} o \texttt{f}
      \item \texttt{1} o \texttt{0}
      \item \texttt{yes} o \texttt{no}
      \item \texttt{NULL}
    \end{itemize}

  \item \textbf{Tipo numerico esatto}:
    Rappresenta numeri interi o con parte decimale fissa.
    \begin{itemize}
      \item Valori interi:
        \begin{itemize}
          \item \texttt{SMALLINT}: 16bit con segno
          \item \texttt{INTEGER}: 32bit con segno
        \end{itemize}

      \item Valori decimali:
        \begin{itemize}
          \item 
\begin{lstlisting}[language=SQL, numbers=none]
NUMERIC | DECIMAL [(precisione [, scala])]
\end{lstlisting}
            \begin{itemize}
              \item \texttt{precisione}: numero totale di cifre (intere + decimali)
              \item \texttt{scala}: numero di cifre decimali
              \item \texttt{|}: indica che è possibile utilizzare sia \texttt{NUMERIC} che \texttt{DECIMAL}
            \end{itemize}
        \end{itemize}
    \end{itemize}

  \item \textbf{Tipo numerico approssimato}:
    Rappresenta numeri in virgola mobile \textbf{approssimati}:
    \begin{itemize}
      \item \texttt{REAL}: 6 cifre decimali
      \item \texttt{DOUBLE PRECISION}: 15 cifre decimali
    \end{itemize}

  \item \textbf{Tipo temporale}:
    Rappresenta istanti di tempo:
    \begin{itemize}
      \item 
\begin{lstlisting}[language=SQL, numbers=none]
DATE
\end{lstlisting}
        Rappresenta una data (fra apici) nel formato ISO: \texttt{YYYY-MM-DD} (anno-mese-giorno)

      \item 
\begin{lstlisting}[language=SQL, numbers=none]
TIME [(precisione)][WITH TIME ZONE]
\end{lstlisting}
      Rappresenta un'orario del tipo \texttt{hh:mm:ss} (ore-minuti-secondi). 
      \begin{itemize}
        \item \texttt{precisione}: numero di cifre decimali per le frazioni di secondo
        \item \texttt{WITH TIME ZONE}: indica che l'orario include anche il fuso orario
          nei seguenti possibili formati:
          \begin{itemize}
            \item \texttt{[+-]hh:mm}: rappresenta la differenza rispetto all'ora di
              Greenwich
            \item Nome del fuso orario (es. \texttt{UTC}, \texttt{CET}, \texttt{EST}, ecc.)
          \end{itemize}
      \end{itemize}

      \item 
\begin{lstlisting}[language=SQL, numbers=none]
TIMESTAMP [(precisione)][WITH TIME ZONE]
\end{lstlisting}
Equivale a \texttt{DATE} + \texttt{TIME}.
    \end{itemize}

  \item \textbf{Tipo intervallo temporale}:
    Rappresenta un intervallo di tempo secondo la seguente sintassi:
\begin{lstlisting}[language=SQL, numbers=none]
INTERVAL PrimaUnitaDiTempo [TO UltimaUnitaDiTempo]
\end{lstlisting}
    dove \texttt{PrimaUnitaDiTempo} e \texttt{UltimaUnitaDiTempo} definiscono le
    unità di misura che devono essere utilizzate, dalla più precisa alla meno precisa.
    \begin{itemize}
      \item \texttt{INTERVAL YEAR}
      \item \texttt{INTERVAL MONTH}
      \item \texttt{INTERVAL DAY}
      \item \texttt{INTERVAL YEAR TO MONTH}
      \item \texttt{INTERVAL DAY TO HOUR}
      \item ecc...
    \end{itemize}
    Ad esempio \texttt{INTERVAL "1 DAY"}
\end{itemize}

\subsubsection{Domini definiti dall'utente}
Per creare un dominio si usa il seguente comando:
\begin{lstlisting}[language=SQL]
CREATE DOMAIN nome AS tipoBase [valoreDefault][vincolo]
\end{lstlisting}
Definisce un dominio utilizzabile in definizioni di relazioni, anche con vincoli
e valori default.
Il dominio può essere definito a partire da un dominio elementare oppure da un altro
dominio definito dall‘utente

\vspace{2em}
\noindent

\begin{example}
  Alcuni esempi di domini definiti dall'utente sono i seguenti:

\begin{lstlisting}[language=SQL]
CREATE DOMAIN Voto                     -- Nome dominio
  AS SMALLINT                          -- Tipo base
  DEFAULT NULL                         -- Valore default
  CHECK ( value >=18 AND value <= 30 ) -- Vincolo
\end{lstlisting}

\begin{lstlisting}[language=SQL]
CREATE DOMAIN giorniSettimana AS CHARACTER (3)
  CHECK (VALUE IN (
    "LUN",
    "MAR",
    "MER",
    "GIO",
    "VEN",
    "SAB",
    "DOM"
  ));
\end{lstlisting}
\vspace{1em}
\noindent

\end{example}

\subsubsection{Vincoli intrarelazionali}
I vincoli intrarelazionali sono vincoli che si applicano a livello di relazione, cioè a
una singola tabella. I principali tipi di vincoli intrarelazionali sono:
\begin{itemize}
  \item \texttt{NOT NULL}: Il valore nullo non è ammesso per un attributo
\begin{lstlisting}[language=SQL, numbers=none]
nome VARCHAR(20) NOT NULL
\end{lstlisting}


  \item \texttt{DEFAULT}: Specifica un valore di default per un attributo quando
    il comando di inserimento non specifica nessun valore per l'attributo

\begin{lstlisting}[language=SQL, numbers=none]
cognome VARCHAR(20) NOT NULL DEFAULT "VUOTO"
\end{lstlisting}

  \item \texttt{UNIQUE}: Impone che i valori di un attributo (o insieme di attributi)
    siano una superchiave, cioè righe diverse della tabella non possono avere gli stessi
    valori

    \begin{itemize}
      \item 
        Il seguente vincolo impone che non ci siano due righe che abbiano sia il nome
        che il cognome uguale
\begin{lstlisting}[language=SQL]
Nome VARCHAR(20),
Cognome VARCHAR(20),
UNIQUE(Nome,Cognome)
\end{lstlisting}

    \item 
      Il seguente vincolo impone che non ci siano due righe che abbiano lo stesso
      nome o lo stesso cognome
\begin{lstlisting}[language=SQL]
Nome VARCHAR(20) UNIQUE,
Cognome VARCHAR(20) UNIQUE
\end{lstlisting}
    \end{itemize}

  \item \texttt{PRIMARY KEY}: Identifica l'attributo che rappresenta la chiave primaria
    della relazione. Impone i seguenti vincoli:
    \begin{itemize}
      \item Si usa una sola volta per tabella
      \item Implica il vincolo \texttt{NOT NULL}
    \end{itemize}
    Si può definire in due modi:
    \begin{itemize}
      \item Nella definizione dell'attributo se la chiave primaria è composta da un
        solo attributo:
\begin{lstlisting}[language=SQL, numbers=none]
matricola CHAR(6) PRIMARY KEY;
\end{lstlisting}

      \item Nel vincolo della tabella se la chiave primaria è composta da più attributi:
\begin{lstlisting}[language=SQL]
nome VARCHAR(20),
cognome VARCHAR(20),
PRIMARY KEY(nome, cognome)
\end{lstlisting}
    \end{itemize}

  \item \texttt{CHECK}: Specifica un vincolo generico che devono soddisfare le tuple
    della tabella. Il vincolo è soddisfatto se la sua espressione è \textbf{vera}
    o \textbf{nulla}.

\begin{lstlisting}[language=SQL, numbers=none]
CREATE TABLE Impiegato (
  matricola CHAR(6) PRIMARY KEY,
  nome VARCHAR(20) NOT NULL,
  cognome VARCHAR(20) NOT NULL,
  qualifica VARCHAR(20),
  stipendio NUMERIC(8,2) DEFAULT 500.00 NOT NULL
  CHECK(stipendio >= 0.0), --check di attributo
  UNIQUE(cognome, nome),
  CHECK(nome <> cognome) --check di tabella
);
\end{lstlisting}
\end{itemize}

\subsubsection{Vincoli interrelazionali}
Un vincolo interrelazionale (o di integrità referenziale) crea un legame tra i valori
di un attributo (o insieme di attributi) \( A \) della tabella corrente, detta
\textbf{interna o slave}, e i valori di un attributo (o di un insieme di attributi)
\( B \) di un'altra tabella detta \textbf{esterna o master}.

Il vincolo impone che:
\begin{itemize}
  \item In ogni tupla della tabella \textbf{interna}, il valore di
    \( A \) (se è diverso dal valore nullo), deve essere presente tra i valori di \( B \) 
    nella tabella esterna.

  \item L'attributo \( B \) della tabella \textbf{esterna} deve essere soggetto ad
    un vincolo \texttt{UNIQUE} o \texttt{PRIMARY KEY}. Può anche non essere la chiave
    primaria, ma deve essere identificante per le tuple della tabella esterna.
\end{itemize}
Per definire un vincolo di integrità referenziale si usa la seguente sintassi:
\begin{itemize}
  \item \texttt{REFERENCES}: Se nel vincolo è coinvolto solo un attributo, si
    specifica la tabella esterna e il relativo attributo coinvolto nel vincolo.

\begin{lstlisting}[language=SQL]
CREATE TABLE Impiegato ( -- Tabella interna
  Matricola CHARACTER(6) PRIMARY KEY,
  Nome CHARACTER VARYING(20) NOT NULL,
  Cognome CHARACTER VARYING(20) NOT NULL,
  Dipart CHARACTER VARYING(15)
    REFERENCES Dipartimento(NomeDip), -- Tabella esterna
  Ufficio NUMERIC(3),
  Stipendio NUMERIC(9) DEFAULT 0,
  UNIQUE (Cognome, Nome)
);
\end{lstlisting}

  \item \texttt{FOREIGN KEY}: Se nel vincolo è coinvolto un insieme di attributi,
    si specifica l'insieme di attributi della tabella interna coinvolti nel vincolo.
    Poi con il vincolo \texttt{REFERENCES} si specifica
    la tabella esterna e l'insieme di attributi della tabella esterna coinvolti nel
    vincolo.

\begin{lstlisting}[language=SQL]
CREATE TABLE Impiegato ( -- Tabella interna
  Matricola CHARACTER(6) PRIMARY KEY,
  Nome CHARACTER VARYING(20) NOT NULL,
  Cognome CHARACTER VARYING(20) NOT NULL,
  Dipart CHARACTER VARYING(15)
    REFERENCES Dipartimento(NomeDip),
  Ufficio NUMERIC(3),
  Stipendio NUMERIC(9) DEFAULT 0,
  FOREIGN KEY (Nome, Cognome)
    REFERENCES Anagrafica (Nome, Cognome) -- Tabella esterna
);
\end{lstlisting}

\end{itemize}

\subsection{CREATE TABLE}
Definisce uno schema di relazione e ne crea un'istanza vuota specificando attributi,
domini e vincoli secondo la seguente sintassi:
\begin{lstlisting}[language=SQL]
CREATE TABLE tabella (
     attributo dominio [ valoreDefault ] { vincolo }
  {, attributo dominio [ valoreDefault ] { vincolo } }
  {, vincoloTabella }
);
\end{lstlisting}
dove:
\begin{itemize}
  \item \texttt{[]}: indicano parti opzionali
  \item \texttt{[ valoreDefault ]}: indica un eventuale valore default
  \item \texttt{\{\}}: indica un elemento che può essere ripetuto zero o più volte
  \item \texttt{\{ vincolo \}}: indica eventuali (zero o più) vincoli di integrità per
    ogni attributo
\end{itemize}

\begin{example}
  Un esempio di creazione di una tabella è il seguente:
\begin{lstlisting}[language=SQL]
CREATE TABLE Impiegato (
  Matricola CHARCTER(6) PRIMARY KEY,
  Nome      CHARCTER    VARYING(20) NOT NULL,
  Cognome   CHARCTER    VARYING(20) NOT NULL,
  Dipart    CHARCTER    VARYING(15),
  Stipendio NUMERIC(9)  DEFAULT 0,

  -- Vincoli di integrita'
  FOREIGN KEY(Dipart)
    REFERENCES Dipartimento(NomeDip), 
  UNIQUE (Cognome, Nome)
)
\end{lstlisting}
\end{example}

\subsection{Modifiche degli schemi}
\subsubsection{DROP DOMAIN}
Permette di rimuovere i domini definiti dall'utente secondo la seguente sintassi:
\begin{lstlisting}[language=SQL, numbers=none]
DROP DOMAIN nomeDominio
\end{lstlisting}

\subsubsection{DROP TABLE}
Permette di rimuovere le tabelle secondo la seguente sintassi:
\begin{lstlisting}[language=SQL, numbers=none]
DROP TABLE nomeTabella
\end{lstlisting}

\subsubsection{ALTER}
Effettua modifiche a domini, tabelle, ecc...
\begin{itemize}
  \item \textbf{Aggiunta di un nuovo attributo} (\texttt{ADD COLUMN}):
    \begin{itemize}
      \item Sintassi:
\begin{lstlisting}[language=SQL, numbers=none]
ALTER TABLE tabella ADD COLUMN attributo tipo;
\end{lstlisting}

      \item Esempio:
\begin{lstlisting}[language=SQL, numbers=none]
ALTER TABLE impiegato ADD COLUMN stipendio NUMERIC(8,2);
\end{lstlisting}
    \end{itemize}

  \item \textbf{Rimozione di un attributo} (\texttt{DROP COLUMN}):
    \begin{itemize}
      \item Sintassi:
\begin{lstlisting}[language=SQL, numbers=none]
ALTER TABLE tabella DROP COLUMN attributo;
\end{lstlisting}

      \item Esempio:
\begin{lstlisting}[language=SQL, numbers=none]
ALTER TABLE impiegato DROP COLUMN stipendio;
\end{lstlisting}
    \end{itemize}

  \item \textbf{Modifica valore di default di un attributo} (\texttt{ALTER COLUMN}):
    \begin{itemize}
      \item Sintassi:
\begin{lstlisting}[language=SQL, numbers=none]
ALTER TABLE tabella ALTER COLUMN attributo { SET DEFAULT valore | DROP DEFAULT };
\end{lstlisting}

      \item Esempio:
\begin{lstlisting}[language=SQL, numbers=none]
ALTER TABLE impiegato ALTER COLUMN stipendio SET DEFAULT 1000.00;
\end{lstlisting}
    \end{itemize}

  \item \textbf{Ridenominazione di un attributo} (\texttt{RENAME COLUMN}):
    \begin{itemize}
      \item Sintassi:
\begin{lstlisting}[language=SQL, numbers=none]
ALTER TABLE tabella RENAME COLUMN vecchioAttributo TO nuovoAttributo;
\end{lstlisting}
      \item Esempio:
\begin{lstlisting}[language=SQL, numbers=none]
ALTER TABLE impiegato RENAME COLUMN stipendio TO salario;
\end{lstlisting}
    \end{itemize}
\end{itemize}

\subsection{Cancellazione dati in una tabella}
Le tuple di una tabella vengono cancellate con il comando \texttt{DELETE} secondo
la seguente sintassi:
\begin{lstlisting}[language=SQL, numbers=none]
DELETE FROM tabella [WHERE condizione];
\end{lstlisting}
\begin{itemize}
  \item \texttt{condizione}: è un'espressione booleana che seleziona quali tuple
    cancellare

  \item Se \texttt{WHERE} non è presente \textbf{tutte le tuple saranno cancellate}
\end{itemize}
Ad esempio:
\begin{lstlisting}[language=SQL, numbers=none]
DELETE FROM impiegato WHERE matricola = "A001";
\end{lstlisting}

\subsection{Inserimento dati in una tabella}
Le tuple di una tabella vengono inserite con il comando \texttt{INSERT} secondo
la seguente sintassi:
\begin{lstlisting}[language=SQL, numbers=none]
INSERT INTO tabella [(attributo1, attributo2, ...)]
VALUES (valore1, valore2, ...);
\end{lstlisting}
oppure
\begin{lstlisting}[language=SQL, numbers=none]
INSERT INTO tabella [(attributo1, attributo2, ...)]
SELECT (...);
\end{lstlisting}
\begin{itemize}
  \item Se la lista degli attributi non è presente, i valori devono essere specificati
    nell'ordine in cui sono stati definiti nella tabella.

  \item Se un attributo non è specificato, viene utilizzato il valore di default se
    è definito, altrimenti viene utilizzato il valore nullo.
\end{itemize}
Ad esempio:
\begin{lstlisting}[language=SQL, numbers=none]
INSERT INTO impiegato (matricola, nome, cognome, dipartimento)
VALUES ("A001", "Mario", "Rossi", "Informatica");
\end{lstlisting}

\subsubsection{Aggiornamento dati in una tabella}
Le tuple di una tabella vengono aggiornate con il comando \texttt{UPDATE} secondo
la seguente sintassi:
\begin{lstlisting}[language=SQL, numbers=none]
UPDATE tabella
SET attributo1 = valore1, attributo2 = valore2, ...
[WHERE condizione];
\end{lstlisting}
\begin{itemize}
  \item \texttt{condizione}: è un'espressione booleana che seleziona quali tuple aggiornare
  \item Se \texttt{WHERE} non è presente \textbf{tutte le tuple saranno aggiornate}
\end{itemize}
Ad esempio:
\begin{lstlisting}[language=SQL, numbers=none]
UPDATE impiegato
SET stipendio = stipendio * 1.1
WHERE dipartimento = "Informatica";
\end{lstlisting}

\subsection{SELECT}
La sintassi dell’istruzione SELECT è la seguente:
\begin{lstlisting}
SELECT [ DISTINCT ]
[ * | expression [[ AS] output_name ] [, ...] ]
[ FROM from_item [, ...] ]
[ WHERE condition ]
[ GROUP BY grouping_element [, ...] ]
[ HAVING condition [, ...] ]
[ { UNION | INTERSECT | EXCEPT } [ DISTINCT ] other_select ]
[ ORDER BY expression [ ASC | DESC | USING operator ]]
\end{lstlisting}
dove:
\begin{itemize}
  \item \texttt{expression}: è un espressione che determina un attributo
  \item \texttt{from\_item}: è un'espressione che determina una sorgente per gli attributi
  \item \texttt{condition}: è un'espressione booleana per selezionare i valori degli attributi
  \item \texttt{grouping\_element}: è un'espressione per poter eseguire operazioni su più
    valori di un attributo
\end{itemize}

\vspace{1em}
\noindent
La struttura dell'istruzione è la seguente:
\begin{lstlisting}
SELECT ListaAttributi
FROM ListaTabelle
[WHERE Condizione]
\end{lstlisting}
\texttt{SELECT}, \texttt{FROM} e \texttt{WHERE} sono chiamate \textbf{clausole}.
Il \texttt{SELECT} estrae dal prodotto cartesiano delle tabelle elencate nella clausola
\texttt{FROM}, le righe che soddisfano la condizione specificata nella clausola \texttt{WHERE}.
Il risultato dell'interrogazione è una tabella con una riga per ogni riga prodotta dalla clausola
\texttt{FROM} ed eventualmente filtrata dalla clausola \texttt{WHERE} (se esiste), le cui
colonne dipendono dalla lista degli attributi nella clausola \texttt{SELECT}.

\begin{example}
  Prendendo ad esempio la seguente base di dati:
  \begin{figure}[H]
    \centering
    \includegraphics[width=0.8\textwidth]{esempio_db}
    \caption{Esempio di base di dati}
  \end{figure}
  \noindent
  Un esempio di istruzione SELECT è il seguente:

  Estrarre lo stipendio degli impiegati di cognome 'Rossi'
  
\begin{lstlisting}
SELECT Stipendio
FROM Impiegato
WHERE Cognome = 'Rossi';
\end{lstlisting}
  Il risultato di questa interrogazione è la seguente tabella:
  \begin{figure}[H]
    \centering
    \includegraphics[width=0.2\textwidth]{esempio_select}
    \caption{Risultato dell'istruzione SELECT}
  \end{figure}
\end{example}

\vspace{1em}
\noindent
Alcune possibili espressioni sono:
\begin{itemize}
  \item \texttt{*}: seleziona tutti gli attributi:
\begin{lstlisting}
SELECT *
FROM Impiegato;
\end{lstlisting}

  \item \texttt{as}: permette di rinominare un attributo:
\begin{lstlisting}
SELECT Stipendio AS Salario
FROM Impiegato;
\end{lstlisting}
\end{itemize}

\subsection{FROM}
\textbf{Clausola FROM}: specifica le tabelle da cui estrarre i dati:
\begin{lstlisting}
FROM from_item [, ...]
\end{lstlisting}
dove \texttt{from\_item} può essere \textbf{una sola} delle seguenti clausole:
\begin{enumerate}
  \item \textbf{table\_name [[AS] alias [(column\_alias [, ...])]}: Se sono presenti due
    o più tabelle, si esegue il prodotto cartesiano tra tutte le tabelle. Se ci sono
    attributi con lo stesso nome su tabelle diverse, nella SELECT questi attributi
    sono identificati anteponendo il nome della tabella e un '.' (es. \texttt{nomeTabella.nomeAttributo}).

  \item \textbf{(other\_select) [AS] nomeRisultato [(column\_alias [,...])]}: è una SELECT
    innestata. Il risultato della SELECT interna è lo schema con nome \texttt{nomeRisultato}
    su cui fare la SELECT corrente.

  \item \textbf{from\_item [NATURAL] join\_type from\_item [ON condition]}: è la clausola JOIN. Metodo
    più sofisticato di 1) che permette di selezionare un sottoinsieme del prodotto cartesiano
    di due o più tabelle. \texttt{join\_type} specifica il tipo di join.
\end{enumerate}

\begin{example}
  Considerando la seguente base di dati:
  \begin{figure}[H]
    \centering
    \includegraphics[width=0.8\textwidth]{esempio_db_from_1}
  \end{figure}
  \begin{figure}[H]
    \centering
    \includegraphics[width=0.8\textwidth]{esempio_db_from_2}
    \caption{Esempio di base di dati per la clausola JOIN}
  \end{figure}
  \noindent
  Estrarre i nomi ed il cognome degli impiegati e le città in cui lavorano:
\begin{lstlisting}
SELECT Impiegato.Nome, Impiegato.Cognome,
       Dipartimento.Citta'
FROM Impiegato, Dipartimento
WHERE Impiegato.Dipart = Dipartimento.Nome
\end{lstlisting}
  Il risultato di questa interrogazione è la seguente tabella:
  \begin{figure}[H]
    \centering
    \includegraphics[width=0.8\textwidth]{esempio_select_from}
    \caption{Risultato dell'istruzione SELECT con clausola FROM}
  \end{figure}
\end{example}

\subsection{WHERE}
\textbf{Clausola WHERE}: specifica la condizione per selezionare le righe da estrarre:
\begin{lstlisting}
WHERE condition
\end{lstlisting}
La clausola \texttt{WHERE} ammette come argomento un’espressione booleana costruita
combinando predicati semplici con gli operatori AND, OR e NOT
\begin{itemize}
  \item \textbf{AND}: seleziona le righe in cui tutti i predicati sono veri
  \item \textbf{OR}: seleziona le righe in cui almeno un predicato è vero
  \item \textbf{NOT}: seleziona le righe in cui il predicato è falso
\end{itemize}

Ciascun predicato semplice usa gli operatori =, <> (\( <> \iff \neq \ \) ), <, >, <= e >= per confrontare
un’espressione costruita a partire dal valore di un attributo con un valore costante,
oppure un’altra espressione.

\begin{example}
  Considerando la seguente base di dati:
  \begin{figure}[H]
    \centering
    \includegraphics[width=0.8\textwidth]{esempio_db_from_1}
  \end{figure}
  \begin{figure}[H]
    \centering
    \includegraphics[width=0.8\textwidth]{esempio_db_from_2}
    \caption{Esempio di base di dati per la clausola JOIN}
  \end{figure}
  \noindent
  Estrarre i nomi ed il cognome degli impiegati che lavorano
  nell’ufficio 20 del dipartimento Amministrazione.
\begin{lstlisting}
SELECT Nome, Cognome
FROM Impiegato
WHERE Uffico = 20 AND
      Dipart = 'Amministrazione'
\end{lstlisting}
  Il risultato di questa interrogazione è la seguente tabella:
  \begin{figure}[H]
    \centering
    \includegraphics[width=0.3\textwidth]{esempio_select_where}
    \caption{Risultato dell'istruzione SELECT con clausola WHERE}
  \end{figure}
\end{example}

\begin{example}
  Considerando la base di dati precedente:

  Estrarre i nomi ed il cognome degli impiegati che lavorano nel dipartimento
  Amministrazione o nel dipartimento Produzione.
\begin{lstlisting}
SELECT Nome, Cognome
FROM Impiegato
WHERE Dipart = 'Amministrazione'
      OR Dipart = 'Produzione'
\end{lstlisting}
  Il risultato di questa interrogazione è la seguente tabella:
  \begin{figure}[H]
    \centering
    \includegraphics[width=0.3\textwidth]{esempio_select_where_or}
    \caption{Risultato dell'istruzione SELECT con clausola WHERE e operatore OR}
  \end{figure}
\end{example}

\begin{example}
  Considerando la base di dati precedente:

  Estrarre i nomi degli impiegati di cognome Rossi che lavorano
  dipartimento Amministrazione o nel dipartimento Produzione.
\begin{lstlisting}
SELECT Nome
FROM Impiegato
WHERE Cognome = 'Rossi' AND
      (Dipart = 'Amministrazione' OR
      Dipart = 'Produzione')
\end{lstlisting}
  Il risultato di questa interrogazione è la seguente tabella:
  \begin{figure}[H]
    \centering
    \includegraphics[width=0.2\textwidth]{esempio_select_where_and_or}
    \caption{Risultato dell'istruzione SELECT con clausola WHERE e operatori AND e OR}
  \end{figure}
\end{example}

\subsubsection{LIKE}
Per il confronto tra stringhe, SQL mette a disposizione anche
l’operatore LIKE che consente verificare se una stringa rispetta un certo
"pattern" (oppure ILIKE per case insentive). I due caratteri speciali per il confronto
sono:
\begin{itemize}
  \item \texttt{\_}: confronto con un carattere arbitrario
  \item \texttt{\%}: confronto con una stringa di lunghezza arbitraria (anche 0) di
    caratteri arbitrari
\end{itemize}
\begin{example}
  Un esempio di pattern è il seguente:
\begin{lstlisting}
_o%i
\end{lstlisting}
  che equivale a una stringa che ha una o in seconda posizione, seguita da una
  sequenza qualsiasi di caratteri (anche vuota) e che termina con una i.
\end{example}

\begin{example}
  Considerando la base di dati precedente:

  Estrarre gli impiegati che hanno un cognome cha ha una "o"
  in seconda posizione e finisce con una "I".

\begin{lstlisting}
SELECT *
FROM Impiegato
WHERE Cognome LIKE '_o%i';
\end{lstlisting}
  Il risultato di questa interrogazione è la seguente tabella:
  \begin{figure}[H]
    \centering
    \includegraphics[width=0.7\textwidth]{esempio_select_where_like}
    \caption{Risultato dell'istruzione SELECT con clausola WHERE e operatore LIKE}
  \end{figure}
\end{example}

\subsubsection{SIMILAR TO}
L’operatore SIMILAR TO è un LIKE più espressivo che accetta,
come pattern, un sottoinsieme delle espressioni regolari POSIX
(versione SQL).

La sintassi delle espressioni regolari è:
\begin{itemize}
  \item \texttt{\_}: 1 carattere qualsiasi
  \item \texttt{\%}: 0 o più caratteri qualsiasi
  \item \texttt{*}: ripetizione del precedente match 0 o più volte
  \item \texttt{+}: ripetizione del precedente match UNA o più volte
  \item \texttt{\{n,m\}}: ripetizione del precedente match almeno n e non più di m volte
  \item \texttt{[...]}: ... è un elenco di caratteri ammissibili
\end{itemize}

\begin{example}
  Considerando la seguente base di dati:
  \begin{figure}[H]
    \centering
    \includegraphics[width=0.8\textwidth]{esempio_db_similar_to}
    \caption{Esempio di base di dati per la clausola SIMILAR TO}
  \end{figure}
  Estrarre gli studenti con cognome che inizia con ’A’ o ’B’, o ’D’, o ’N’ e finisce con ’a’:
\begin{lstlisting}
SELECT cognome, nome, citta' FROM Studente
WHERE cognome SIMILAR TO '[ABDN]{1}%a';
\end{lstlisting}
  Il risultato di questa interrogazione è evidenziato in giallo.
\end{example}

\subsubsection{BETWEEN}
L’operatore BETWEEN consente di verificare se un valore è compreso tra due valori
estremi (inclusi). La sintassi è la seguente:
\begin{lstlisting}
expression BETWEEN value1 AND value2
\end{lstlisting}

\begin{example}
  Considerando la seguente base di dati:
  \begin{figure}[H]
    \centering
    \includegraphics[width=0.8\textwidth]{esempio_db_between}
    \caption{Esempio di base di dati per la clausola BETWEEN}
  \end{figure}
  Estrarre tutti gli studenti che hanno matricola tra ’IN0002’ e ’IN0004’.
\begin{lstlisting}
SELECT cognome, nome, matricola FROM Studente
WHERE matricola BETWEEN 'IN0002' AND 'IN0004';
\end{lstlisting} 
  Il risultato di questa interrogazione è il seguente:
  \begin{figure}[H]
    \centering
    \includegraphics[width=0.3\textwidth]{esempio_select_where_between}
    \caption{Risultato dell'istruzione SELECT con clausola WHERE e operatore BETWEEN}
  \end{figure}
\end{example}

\subsubsection{IN}
Nella clausola WHERE può apparire l’operatore [NOT] IN per
testare l’appartenenza di un valore ad un insieme. La sintassi è la seguente:
\begin{lstlisting}
WHERE attributo [NOT] IN ( valore [, ...] );
\end{lstlisting}

\subsubsection{IS NULL}
Per verificare se un attributo ha un valore nullo o meno si utilizza il predicato
IS NULL che ritorna true se l'attributo ha valore nullo. Oppure IS NOT NULL.

\begin{lstlisting}
WHERE attributo IS [ NOT ] NULL;
\end{lstlisting}

SQL-2 implementa la logica a 3 valori, cioè un predicato semplice restituisce il valore
NULL quando uno qualsiasi dei termini del predicato ha valore nullo. Il predicato IS NULL
è una eccezione, restituendo sempre il valore vero oppure false, mai il valore NULL.

\subsection{Gestione dei duplicati}
Siccome la rimozione delle righe duplicate è un'operazione costosa, SQL le mantiene
di default. Per togliere i duplicati si può utilizzare la parola chiave DISTINCT dopo SELECT:
\begin{lstlisting}
SELECT DISTINCT Stipendio
FROM Impiegato
WHERE Cognome = 'Rossi';
\end{lstlisting}

Il costrutto ALL è invece usato per mantenere i duplicati, ma è superfluo.

\begin{example}
  Considerando la seguente base di dati:
  \begin{figure}[H]
    \centering
    \includegraphics[width=0.8\textwidth]{esempio_db_distinct}
    \caption{Esempio di base di dati per la clausola DISTINCT}
  \end{figure}
  Estrarre le città delle persone il cui cognome è Rossi
\begin{lstlisting}[language=SQL]
SELECT Citta
  FROM Persona
  WHERE Cognome = 'Rossi'
\end{lstlisting}
  L'output di questa interrogazione è la seguente tabella:
  \begin{table}[H]
    \centering
    \begin{tabular}{|c|}
      \hline
      Città \\
      \hline
      Verona \\
      Verona \\
      Milano \\
      \hline
    \end{tabular}
  \end{table}
  \noindent
  Per togliere i duplicati si può utilizzare la parola chiave DISTINCT dopo SELECT:
\begin{lstlisting}[language=SQL]
SELECT DISTINCT Citta
  FROM Persona
  WHERE Cognome = 'Rossi'
\end{lstlisting}
  L'output di questa interrogazione è la seguente tabella:
  \begin{table}[H]
    \centering
    \begin{tabular}{|c|}
      \hline
      Città \\
      \hline
      Verona \\
      Milano \\
      \hline
    \end{tabular}
  \end{table}
\end{example}

\subsection{JOIN}
Le interrogazioni con il join possono essere definite in due modi:
\begin{enumerate}
  \item Senza l’uso di nuova sintassi, e specifica il legame
    tra le tabelle coinvolte a livello della clausola WHERE.

  \item Con una sintassi riservata a livello di clausola FROM
\end{enumerate}

% • La clausola FROM ammette come argomento:
% table_name [[AS] name]] [, ...]
% • Si è visto che se sono presenti due o più nomi di tabelle, si esegue
% il prodotto cartesiano tra tutte le tabelle e lo schema del risultato
% può contenere tutti gli attributi del prodotto cartesiano.
% • Il prodotto cartesiano di due o più tabelle è un CROSS JOIN.
% • A partire da SQL-2, esistono altri tipi di JOIN (join_type):
% INNER JOIN, LEFT [OUTER] JOIN, RIGHT [OUTER] JOIN e FULL
% [OUTER] JOIN.
La clausola FROM ammette come argomento:
\begin{lstlisting}
table_name [[AS] name]] [, ...]
\end{lstlisting}
Si è visto che se sono presenti due o più nomi di tabelle, si esegue il prodotto
cartesiano tra tutte le tabelle e lo schema del risultato può contenere tutti gli attributi
del prodotto cartesiano. Il prodotto cartesiano di due o più tabelle è un CROSS JOIN.

A partire da SQL-2, esistono altri tipi di JOIN (join\_type):
\begin{itemize}
  \item \textbf{INNER JOIN}: restituisce solo le righe che soddisfano la condizione di join
  \item \textbf{LEFT [OUTER] JOIN}: restituisce tutte le righe della tabella di sinistra e le righe della tabella di destra che soddisfano la condizione di join, e NULL per le righe della tabella di destra che non soddisfano la condizione di join
  \item \textbf{RIGHT [OUTER] JOIN}: restituisce tutte le righe della tabella di destra e le righe della tabella di sinistra che soddisfano la condizione di join, e NULL per le righe della tabella di sinistra che non soddisfano la condizione di join
  \item \textbf{FULL [OUTER] JOIN}: restituisce tutte le righe di entrambe le tabelle, e NULL per le righe che non soddisfano la condizione di join
\end{itemize}
La sintassi generale è:
\begin{lstlisting}
table_name [ NATURAL ] join_type table_name [ON join_condition [, ...]].
\end{lstlisting}
dove \texttt{join\_condition} è un’espressione booleana che
seleziona le tuple del join da aggiungere al risultato. Le
tuple selezionate possono essere poi filtrate con la
condizione della clausola WHERE

\begin{example}[Prodotto Cartesiano]
  \begin{figure}[H]
    \centering
    \includegraphics[width=0.8\textwidth]{esempio_db_join}
    \caption{Esempio di base di dati per la clausola JOIN}
  \end{figure}
\end{example}

\begin{example}
  Considerando la seguente base di dati:
  \begin{figure}[H]
    \centering
    \includegraphics[width=0.8\textwidth]{esempio_db_join_2}
    \caption{Esempio di base di dati per la clausola JOIN}
  \end{figure}

\begin{lstlisting}[language=SQL]
SELECT i.Nome, i.Cognome, d.Citta
  FROM Impiegato as i, Dipartimento as d -- Prodotto cartesiano
  WHERE i.Dipart = d.Nome
\end{lstlisting}
\end{example}

\noindent
Il JOIN si può spostare nella clausola FROM, a livello delle tabelle (questa sintassi
sarà utilizzata per gli esercizi del corso). La sintassi è la seguente:
\begin{lstlisting}[language=SQL]
SELECT ListaAttributi
  FROM Tabella1 JOIN Tabella2 ON ( CondizioneDiJoin )
  [WHERE AltraCondizione]
\end{lstlisting}

\begin{example}
  Ad esempio:
\begin{lstlisting}[language=SQL]
SELECT i.Nome, d.Citta
FROM Impiegato AS i JOIN Dipartimento as d ON ( i.Dipart = d.Nome )
\end{lstlisting}
\end{example}

\begin{example}
  Considerando la seguente base di dati:
  \begin{figure}[H]
    \centering
    \includegraphics[width=0.8\textwidth]{esempio_db_distinct}
    \caption{Esempio di base di dati per la clausola DISTINCT}
  \end{figure}
  Selezionare i padri di persone che guadagnano più di 20
\begin{lstlisting}[language=SQL]
SELECT DISTINCT Padre
  FROM Persone JOIN Paternita ON ( Figlio = Nome )
  WHERE Reddito > '20' 
\end{lstlisting}
\end{example}

\subsubsection{INNER JOIN}
La sintassi è la seguente:
\begin{lstlisting}
table1 [ INNER ] JOIN table2 ON join_condition [, ...]
\end{lstlisting}
Rappresenta il tradizionale \( \Theta  \) join dell’algebra relazionale. Combina ciascuna
riga \( r_1 \) di table1 con ciascuna riga di table2 che soddisfa la condizione della clausola
ON.

\begin{example}
  \begin{figure}[H]
    \centering
    \includegraphics[width=0.8\textwidth]{esempio_db_inner_join}
    \caption{Esempio di base di dati per la clausola INNER JOIN}
  \end{figure}
\end{example}

\subsubsection{LEFT OUTER JOIN}
La sintassi è la seguente:
\begin{lstlisting}
table1 LEFT [ OUTER ] JOIN table2 ON join_condition [, ...].
\end{lstlisting}
Si esegue un INNER JOIN. Poi, per ciascuna riga \( r_1 \) di table1 che non soddisfa la
condizione con qualsiasi riga di table2, si aggiunge una riga al risultato con i valori
di \( r_1 \) e assegnando NULL agli altri attributi.

\begin{example}
  \begin{figure}[H]
    \centering
    \includegraphics[width=0.8\textwidth]{esempio_db_left_outer_join}
    \caption{Esempio di base di dati per la clausola LEFT OUTER JOIN}
  \end{figure}
\end{example}

\textbf{ATTENZIONE}: Il LEFT [OUTER] JOIN non è simmetrico, cioè con le stesse tabelle
si possono ottenere risultati diversi a seconda dell’ordine in cui si scrivono le tabelle.

\subsubsection{RIGHT OUTER JOIN}
La sintassi è la seguente:
\begin{lstlisting}
table1 RIGHT [ OUTER ] JOIN table2 ON join_condition [, ...].
\end{lstlisting}
Si esegue un INNER JOIN. Poi, per ciascuna riga \( r_2 \) di table2 che non soddisfa la
condizione con qualsiasi riga di table1, si aggiunge una riga al risultato con i valori
di \( r_2 \) e assegnando NULL agli altri attributi.

\begin{example}
  \begin{figure}[H]
    \centering
    \includegraphics[width=0.8\textwidth]{esempio_db_right_outer_join}
    \caption{Esempio di base di dati per la clausola RIGHT OUTER JOIN}
  \end{figure}
\end{example}

\subsubsection{FULL OUTER JOIN}
La sintassi è la seguente:
\begin{lstlisting}
table1 FULL [ OUTER ] JOIN table2 ON join_condition [, ...].
\end{lstlisting}
È equivalente a:
\begin{lstlisting}
INNER JOIN + LEFT OUTER JOIN + RIGHT OUTER JOIN
\end{lstlisting}
Non è equivalente a CROSS JOIN!

\subsection{Uso di variabili (alias)}
SQL permette di associare un nome alternativo (alias) alle tabelle
che compaiono nella clausola FROM e viene usato per far riferimento alla tabella nel
contesto dell’interrogazione.

\begin{example}
  Considerando la seguente base di dati:
  \begin{figure}[H]
    \centering
    \includegraphics[width=0.8\textwidth]{esempio_db_alias}
    \caption{Esempio di base di dati per l'uso di alias}
  \end{figure}
  Estrarre tutti gli impiegati che hanno lo stesso cognome ma nome
  diverso, del dipartimento Produzione.
\begin{lstlisting}[language=SQL]
SELECT i1.Cognome, I1.Nome
  FROM Impiegato i1, Impiegato i2
  WHERE i1.Cognome = i2.Cognome AND
  i1.Nome <> i2.Cognome AND
  i2.Dipart = 'Produzione'
\end{lstlisting}
  Questa interrogazione confronta ciascuna riga di Impiegato con tutte le righe
  di Impiegato associate al dipartimento Produzione.
  Si può immaginare che al momento della definizione dell’alias vengano
  create due copie della tabella Impiegato, ciascuna associata ad una
  delle due variabili i1 e i2. Le due copie contengono tutte le
  righe di Impiegato Sulla copia i2 vengono selezionate
  solo le righe del dipartimento 'Produzione'.
  \begin{figure}[H]
    \centering
    \includegraphics[width=0.8\textwidth]{rappresentazione_alias}
    \caption{Esempio di base di dati per l'uso di alias e JOIN}
  \end{figure}
\end{example}

\subsection{Ordinamento del risultato}
Una relazione è un insieme di tuple non ordinato, ma spesso c'è bisogno di costruire
un ordine sulle righe. Per questo SQL mette a disposizione la clausola ORDER BY che permette
di ordinare le righe del risultato in base ai valori di uno o più attributi. La sintassi è la
seguente:
\begin{lstlisting}
ORDER BY attributo [{ ASC | DESC }] [, ... ];
\end{lstlisting}
Le righe vengono ordinate per il primo attributo nell’elenco, a parità di valore
del primo attributo si usa il secondo, e così via in sequenza. Di default l’ordinamento è ASC.

\begin{example}
  Considerando la seguente base di dati:
  \begin{figure}[H]
    \centering
    \includegraphics[width=0.8\textwidth]{esempio_db_order_by}
    \caption{Esempio di base di dati per la clausola ORDER BY}
  \end{figure}
  Estrarre il nome, l’età ed il reddito delle persone con meno
  di trenta anni in ordine alfabetico.
\begin{lstlisting}[language=SQL]
SELECT Nome, Eta, Reddito
  FROM Persone
  WHERE Eta < '30'
  ORDER BY Eta
\end{lstlisting}
\end{example}

\begin{example}
  Considerando la seguente base di dati:
  \begin{figure}[H]
    \centering
    \includegraphics[width=0.8\textwidth]{esempio_db_order_by_2}
    \caption{Esempio di base di dati per la clausola ORDER BY}
  \end{figure}
  Estrarre il contenuto della tabella Automobile ordinato in
  base alla marca (in modo discendente) e al modello.
\begin{lstlisting}[language=SQL]
SELECT *
  FROM Automobile
  ORDER BY Marca DESC, Modello
\end{lstlisting}
\end{example}

\subsection{Operatori insiemistici}
Gli operatori insiemistici disponibili sono:
\begin{lstlisting}
SELECT
  { UNION | INTERSECT| EXCEPT
  [ALL] SELECT }
\end{lstlisting}
Gli operatori insiemistici, al contrario del resto del linguaggio, assumono come default
di eseguire un’eliminazione dei duplicati, a meno che si usi ALL. Sono sufficienti domini
compatibili ed ugual numero di attributi. La corrispondenza degli attributi si basa sulla
posizione, e se i nomi degli attributi sono diversi di solito il sistema usa il nome del
primo operando.

\subsection{Operatori aggregati}
Gli operatori aggregati permettono di eseguire operazioni su più righe per ottenere un unico
valore. Gli operatori aggregati disponibili sono:
\begin{itemize}
  \item \textbf{COUNT}
  \item \textbf{SUM}
  \item \textbf{AVG}
  \item \textbf{MIN}
  \item \textbf{MAX}
\end{itemize}
Tutte le condizioni specificate vengono valutate su una tupla alla volta e gli operatori
permettono di valutare proprietà che dipendono da insiemi di tuple.

Prima viene eseguita l’interrogazione considerando solo le clausole
FROM e WHERE, quindi l’operatore di aggregazione viene applicato alla
tabella contenente il risultato.

\subsubsection{COUNT}
Conta il numero di righe di una tabella o il numero di valori di uno o più attributi.
La sintassi è la seguente:
\begin{lstlisting}
COUNT ( <* | [DISTNCT | ALL] ListaAttributi> )
\end{lstlisting}
dove:
% DISTINCT ListaAttributi: conta il numero di valori diversi degli
% attributi ListaAttributi
% • ALL ListaAttributi: conta il numero di righe che hanno valori diversi
% da NULL negli attributi in ListaAttributi
% • Comportamento di default
\begin{itemize}
  \item \texttt{*}: conta il numero di righe della tabella

  \item \texttt{DISTINCT ListaAttributi}: conta il numero di valori diversi degli
    attributi ListaAttributi

  \item \texttt{ALL ListaAttributi}: conta il numero di righe che hanno valori diversi da
    NULL negli attributi in ListaAttributi
\end{itemize}

\subsubsection{SUM}
La sintassi è la seguente:
\begin{lstlisting}
SUM ( [DISTINCT | ALL] AttrExpr )
\end{lstlisting}
Restituisce la somma dei valori posseduti dall'espressione.
Ammette come argomento solo espressioni che rappresentano valori numerici o intervalli
di tempo.

\begin{itemize}
  \item DISTINCT elimina i duplicati
  \item ALL trascura solo i valori nulli
\end{itemize}

\subsubsection{MAX/MIN}
La sintassi è la seguente:
\begin{lstlisting}
<MAX | MIN> ( [DISTINCT | ALL] AttrExpr )
\end{lstlisting}
Restituisce il massimo/minimo dei valori posseduti dall’espressione.
Ammettono come argomento solo espressioni su cui è definito un ordinamento (quindi anche
stringhe)

\subsubsection{AVG}
La sintassi è la seguente:
\begin{lstlisting}
AVG ( [DISTINCT | ALL] AttrExpr )
\end{lstlisting}
Restituisce la media dei valori posseduti dall’espressione.

\begin{itemize}
  \item DISTINCT elimina i duplicati
  \item ALL trascura solo i valori nulli
\end{itemize}


\subsection{Interrogazioni con raggruppamento}
Gli operatori di aggregazione possono essere applicati separatamente a sottoinsiemi di righe.
La clausola GROUP BY permette di specificare come dividere le tabelle in sottoinsiemi e
permette di specificare un insieme di attributi secondo cui raggruppare le righe che possiedono
gli stessi valori per questo insieme di attributi.

\vspace{1em}
\noindent
In una interrogazione che fa uso della clausola GROUP BY può comparire
come argomento della SELECT solo un sottoinsieme degli attributi usati nella
clausola GROUP BY oppure operatori di aggregazione.
Ogni sottoinsieme di righe generato dal GROUP BY deve corrispondere ad un’unica riga nella
tabella risultato.

\subsubsection{HAVING}
È possibile applicare delle selezioni ai gruppi generati tramite la clausola GROUP BY
tramite la clausola HAVING. Questa clausola permette di specificare delle condizioni che
devono essere soddisfatte dai gruppi generati dopo l’esecuzione
del raggruppamento tramite GROUP BY.

\vspace{1em}
\noindent
\textbf{Osservazione}:
\begin{itemize}
  \item La clausola WHERE applica delle selezioni su SINGOLE RIGHE, prima dell’eventuale
    raggruppamento tramite GROUP BY.

  \item La clausola HAVING applica delle selezioni su GRUPPI, generati a seguito dell’esecuzione del
    raggruppamento tramite GROUP BY.
\end{itemize}

\subsection{Interrogazioni nidificate}
Un'interrogazione nidificata è un'interrogazione che contiene al suo interno altre
interrogazioni. La nidificazione può avvenire sia nella clausola SELECT, che nella
clausola FROM, che nella clausola WHERE. Il caso tipico è la nidificazione nella clausola
WHERE, dove un predicato contiene un confronto tra un valore (espressione su una singola
riga) ed il risultato di un’interrogazione SQL.

\subsubsection{WHERE}
Se in un predicato si confronta un attributo o una espressione che produce un valore
singolo con il risultato di un’interrogazione, sorge il problema di disomogeneità dei termini
del confronto. I normali operatori di confronto (=, <>, <, >, $\le$ , $\ge$) possono
essere estesi con le parole chiave ANY e ALL:
\begin{itemize}
  \item \textbf{ANY}: la riga soddisfa la condizione se risulta vero il confronto
    (con l’operatore specificato) tra il valore dell’attributo (o espressione) e ALMENO uno
    degli elementi restituiti dall’interrogazione.

  \item \textbf{ALL}: la riga soddisfa la condizione se risulta vero il confronto
    (con l’operatore specificato) tra il valore dell’attributo (o espression) e TUTTI gli elementi restituiti dall’interrogazione.
\end{itemize}
Per il controllo di appartenenza o esclusione rispetto ad un insieme si possono usare anche
gli operatori IN (= ANY) e NOT IN (<> ALL).
L'interrogazione può essere eseguita una sola volta ed il risultato ottenuto può
essere salvato in una tabella temporanea usata per eseguire
ciascun confronto con le righe dalla tabella esterna, però  l’interrogazione nidificata
fa riferimento al contesto dell’interrogazione che la racchiude attraverso l’uso di una
variabile (\textbf{passaggio di binding}), quindi l'interrogazione nidificata
dovrà essere \textbf{valutata separatamente per ciascuna riga} prodotta dalla
valutazione della query esterna.

\subsubsection{Binding e EXISTS}
L'operatore EXISTS riceve come parametro un'interrogazione nidificata e restituisce il
valore vero solo se l'interrogazione fornisce un risultato non vuoto. Si può utilizzare
soltanto quando c'è un passaggio di binding tra la query esterna e quella interna,
cioè quando la query interna fa riferimento ad una variabile che assume valori dalle
righe prodotte dalla query esterna.
\begin{example}
  Estrarre le persone che hanno degli omonimi (ovvero
  persone con lo stesso nome e cognome, ma codice fiscale
  diverso)
\begin{lstlisting}[language=SQL]
SELECT *
  FROM Persona P
  WHERE EXISTS (
    SELECT *
    FROM Persona P1
    WHERE P1.Nome = P.Nome AND
    P1.Cognome = P.Cognome AND
    P1.CodFiscale <> P.CodFiscale
  );
\end{lstlisting}
\end{example}
\begin{example}
  Estrarre le persone che non hanno degli omonimi (ovvero
  persone con lo stesso nome e cognome, ma codice fiscale
  diverso)
\begin{lstlisting}[language=SQL]
SELECT *
  FROM Persona P
  WHERE NOT EXISTS (
    SELECT *
    FROM Persona P1
    WHERE P1.Nome = P.Nome AND
    P1.Cognome = P.Cognome AND
    P1.CodFiscale <> P.CodFiscale
  );
\end{lstlisting}
  oppure
\begin{lstlisting}[language=SQL]
SELECT *
  FROM Persona P
  WHERE NOT IN (
    SELECT Nome, Cognome
    FROM Persona P1
    WHERE P1.CodFiscale <> P.CodFiscale
  );
\end{lstlisting}
\end{example}

\subsubsection{Interrogazioni nidificate nella clausola SELECT}
Le interrogazioni nidificate possono essere specificate anche come
elemento della clausola SELECT.
È necessario che l’interrogazione nidificata restituisca una
sola tupla (riferimenti tramite variabile alla interrogazione esterna).
Spesso le interrogazioni nidificate nel SELECT
sostituiscono l'uso del JOIN.
\begin{example}
  Estrarre per ogni cantante il numero di canzoni di cui è autore.
\begin{lstlisting}[language=SQL]
SELECT DISTINCT Nome,
  NumCanzoni =
    (
      SELECT COUNT(*)
      FROM Autore A
      WHERE A.Nome = C.Nome
    )
  FROM Cantante C
\end{lstlisting}
  equivale a 
\begin{lstlisting}[language=SQL]
  SELECT DISTINCT C.Nome,
    COUNT(*) as NumCanzoni
    FROM Cantante C JOIN Autore A
    ON C.Nome = A.Nome
    GROUP BY C.Nome
\end{lstlisting}
\end{example}

\subsubsection{Interrogazioni nidificate nella clausola FROM}
Le interrogazioni nidificate possono essere utilizzate nella clausola FROM come
sorgente dati su cui eseguire l’interrogazione.
È possibile creare una catena di interrogazioni, dove ciascuna
interrogazione utilizza il risultato della precedente come una
tabella di base.

\begin{example}
  Estrarre per ogni cantante che è anche autore di canzoni il numero
  di canzoni di cui è rispettivamente interprete e autore.
\begin{lstlisting}[language=SQL]
SELECT DISTINCT C.Nome, C.NumCantate, A.NumScritte
  FROM
  ( SELECT Nome, Count(*) as NumCantate FROM Cantante GROUP BY Nome) C
  JOIN
  ( SELECT Nome, Count(*) as NumScritte FROM Autore GROUP BY Nome) A
  ON C.Nome = A.Nome
\end{lstlisting}
  Prima vengono eseguite le interrogazioni nella clausola FROM e i
  risultati di queste vengono usati come normali tabelle.
\end{example}

\subsubsection{Esercizi}

\begin{exercise}
  Considerando la seguente base di dati:
  \begin{itemize}
    \item Turista(\underline{username}, cognome, nome, data\_nascita, nazionalita)

    \item Attrazione(\underline{codice}, nome, descrizione, latitudine, longitudine,
      indirizzo, citta, prov)

    \item Recensione(\underline{attrazione, turista, timestamp}, titolo, testo, voto)
  \end{itemize}
  \noindent
  Con i seguenti vincoli di integrità:
  \begin{itemize}
    \item Recensione.attrazione $\to$ Attrazione
    \item Recensione.turista $\to$ Turista
  \end{itemize}
  \noindent
  Scrivere le seguenti interrogazioni SQL:
  \begin{enumerate}
    \item Trovare tutte le attrazioni della città di Verona il cui nome inizia con la
      lettera f maiuscola o minuscola
\begin{lstlisting}[language=SQL]
SELECT *
  FROM Attrazione
  WHERE nome ILIKE 'f%';
\end{lstlisting}
      oppure
\begin{lstlisting}[language=SQL]
SELECT *
  FROM Attrazione
  WHERE nome LIKE 'f%' OR nome LIKE 'F%';
\end{lstlisting}
      oppure
\begin{lstlisting}[language=SQL]
SELECT *
  FROM Attrazione
  WHERE LOWER(nome) LIKE 'f%'; -- Converte nome in minuscolo;
\end{lstlisting}


    \item Trovare tutti i turisti che nel 2025 \textbf{non} hanno mai rilasciato
      recensioni con un voto maggiore di 2 per attrazioni fuori dalla provincia
      di Verona.
\begin{lstlisting}[language=SQL]
SELECT turista
  FROM Turista T
  EXCEPT -- Le due query devono avere schemi compatibili
  (
    SELECT turista
      FROM Recensione as R
      JOIN Attrazione A
        ON (A.codice = R.attrazione)
      WHERE EXTRACT(YEAR FROM R.timestamp) = '2025'
        AND R.voto > 2 
        AND A.provincia <> 'Verona'
  );
\end{lstlisting}
      Questa query è difficile da cambiare, ad esempio aggiungendo attributi perchè
      bisognerebbe cambiare sia la query esterna che quella interna.
      Un'alternativa è la seguente:
\begin{lstlisting}[language=SQL]
SELECT T.username, T.nome, T.cognome
  FROM Turista as T
  WHERE T.username
  NOT IN
  (
    SELECT turista
      FROM Recensione as R
      JOIN Attrazione A
        ON (A.codice = R.attrazione)
      WHERE EXTRACT(YEAR FROM R.timestamp) = '2025'
        AND R.voto > 2 
        AND A.provincia <> 'Verona'
  );
\end{lstlisting}
      Questa query non richiede che gli schemi siano compatibili, e permette di
      aggiungere facilmente altri attributi alla query esterna.

    \item Trovare tutte le attrazioni che hanno ricevuto nel 2025 un numero di
      recensioni superiore a quelle ricevute nel 2024.
\begin{lstlisting}[language=SQL]
-- Attrazioni con numero recensioni nel 2025
SELECT R1.attrazione, COUNT(*)
  FROM Recensione as R1
  WHERE EXTRACT(YEAR FROM R1.timestamp) = 2025
  GROUP BY R1.attrazione;
  HAVING COUNT(*) > (
    -- Numero recensioni nel 2024
    SELECT COUNT(*)
      FROM Recensione as R2
      WHERE R1.attrazione = R2.attrazione
        AND EXTRACT(YEAR FROM R2.timestamp) = 2024;
  );
\end{lstlisting}
      Nella query annidata si può omettere il GROUP BY perchè grazie al data binding
      si sta già confrontando le recensioni della stessa attrazione.

    \item Trovare la media dei voti massimi rilasciati dagli utenti.
\begin{lstlisting}[language=SQL]
SELECT AVG(R1.voto_max)
  FROM (
    SELECT MAX(R2.voto) as voto_max
      FROM Recensione as R2
      GROUP BY R2.turista;
  ) as R1;
\end{lstlisting}
      Non si può scrivere AVG(MAX()) perchè l'operatore aggregato esterno non riceve
      una lista di valori, ma un singolo attributo che rappresenta il risultato
      dell'operatore aggregato interno.

  \end{enumerate}
\end{exercise}


\begin{exercise}
  \label{exercise:parcheggi}
  Considerando la seguente base di dati:
  \begin{itemize}
    \item Auto(\underline{targa}, num\_telaio, modello, colore, km, cilindrata, posti,
      proprietario)

    \item Parcheggio(\underline{numero}, comune, indirizzo, coperto)

    \item Prenotazione(\underline{codice}, parcheggio, posto, auto, data, ora\_inizio,\\
      ora\_fine, tipo, pagato)
  \end{itemize}
  \noindent
  Con i seguenti vincoli di integrità:
  \begin{itemize}
    \item Prenotazione.parcheggio $\to$ Parcheggio
    \item Prenotazione.auto $\to$ Auto
  \end{itemize}
  \noindent
  Scrivere le seguenti interrogazioni SQL:
  \begin{enumerate}
    \item Trovare tutti i parcheggi coperti di Verona
\begin{lstlisting}[language=SQL]
SELECT *
  FROM Parcheggio
  WHERE 
    comune = 'Verona';
    AND coperto
\end{lstlisting}

    \item Trovare tutti i parcheggi in cui \textbf{non} sono mai state fatte
      prenotazioni online nel 2024. Per ciascun parcheggio riportare il nome
      e il numero di prenotazioni complessive.
\begin{lstlisting}[language=SQL]
-- Tutti i parcheggi
SELECT P.nome, COUNT(*) as num_prenotazionu
  FROM Parcheggio as P
    JOIN Prenotazione as Pr
    ON (P.codice = Pr.parcheggio)
  WHERE P.codice
    NOT IN (
      -- Parcheggi con almeno una prenotazione online nel 2024
      SELECT P2.parcheggio
      FROM  Prenotazione as Pr2
      WHERE Pr2.tipo = 'online'
        AND EXTRACT(YEAR FROM data) = '2024'
    )
  GROUP BY P.nome;
\end{lstlisting}

    \item Trovare tutti i parcheggi in cui vengono effettuate più prenotazioni online
      rispetto a quelle fisihce.
\begin{lstlisting}[language=SQL]
SELECT Pr.parcheggio
  FROM Prenotazione as Pr
  WHERE Pr.tipo = 'online'
  GROUP BY Pr.parcheggio
    HAVING COUNT(*) > (
      SELECT COUNT(*)
        FROM Prenotazione as Pr2
        WHERE Pr2.tipo = 'fisica'
          AND Pr2.parcheggio = Pr.parcheggio -- Data binding
    );
\end{lstlisting}
    Siccome c'è un data binding tra la query esterna e quella interna, la query interna
    viene valutata separatamente per ciascun parcheggio restituito dalla query esterna.

  \end{enumerate}
\end{exercise}

\subsection{Viste}
Le viste permettono di salvare una query che genera una tabella temporanea.
La sintassi è la seguente:
\begin{lstlisting}[language=SQL, numbers=none]
CREATE VIEW nome_vista AS ( query );
\end{lstlisting}
Si da un nome alla query (nome\_vista) che si può usare come se fosse il nome di una
tabella.

\begin{example}
  Considerando la base di dati dell'esercizio \ref{exercise:parcheggi},
  trovare quanti parcheggi del comune di 'Verona' hanno avuto un numero di prenotazioni
  nel 2024 superiore al numero medio di prenotazioni complessive nel 2024 dei parcheggi
  di 'Padova'.

  I passi da seguire sono:
  \begin{enumerate}
    \item Calcolare il numero di prenotazioni del 2024 per parcheggio, per ogni
      parcheggio di 'Verona'

    \item \label{step:padova}
      Calcolare il numero di prenotazioni del 2024 per parcheggio, per ogni
      parcheggio di 'Padova'

    \item Calcolare la media su \ref{step:padova}
  \end{enumerate}

\begin{lstlisting}[language=SQL]
-- Per ogni parcheggio si ottiene:
-- * codice
-- * comune
-- * num prenotazioni
CREATE VIEW NumPrenotazioni as (
  SELECT Pr.codice as parcheggio, P.comune, COUNT(*) as num_pr
  FROM Parcheggio as P
    JOIN Prenotazione as Pr
    ON (P.codice = Pr.parcheggio)
  WHERE EXTRACT(YEAR FROM Pr.data) = '2024'
  GROUP BY Pr.codice, P.comune
);

-- Calcola il numero medio di prenotazioni dei parcheggi
-- di Padova nel 2024
SELECT NP1.parcheggio
FROM NumPrenotazioni as NP1
WHERE NP1.num_pr > (
    SELECT AVG(NP2.num_pr) -- Ritorna 1 valore
    FROM NumPrenotazioni as NP2
    WHERE NP2.comune = 'Padova'
  )
  AND NP1.comune = 'Verona';
  -- Per ogni parcheggio di 'Verona' guardo quali hanno il numero
  -- di prenotazioni > della media di quelli di 'Padova'
\end{lstlisting}
\end{example}

\section{Strutture fisiche e strutture di accesso ai dati}
Le basi di dati risiedono in \textbf{memoria secondaria}, in quanto sono:
\begin{itemize}
  \item \textbf{Grandi}: non possono essere contenute completamente in memoria centrale
  \item \textbf{Persistente}: hanno un tempo di vita che non è limitato all'esecuzione
    dei programmi che le usano
\end{itemize}
Le caratteristiche della memoria secondaria sono:
\begin{itemize}
  \item Non è direttamente utilizzabile dai programmi

  \item I dati sono organizzati \textbf{a blocchi}. La grandezza dei
    blocchi dipende dal sistema operativo.

  \item Le uniche operazioni possibili sono:
    \begin{itemize}
      \item Lettura di un intero blocco
      \item Scrittura di un intero blocco
    \end{itemize}

  \item Il costo di tali operazioni è ordini di grandezza maggiore del costo per
    accedere ai dati in memoria centrale
\end{itemize}

\subsection{Gestore dei buffer}
L'interazione tra memoria secondaria e memoria centrale avviene tramite il
\textbf{gestore del buffer} e consiste nel caricare uno o più \textbf{blocchi
dalla memoria secondaria} in \textbf{pagine} dedicate della \textbf{memoria centrale}
detta \textbf{buffer}. La dimensione delle pagine in memoria centrale sono solitamente
un multiplo della dimensione dei blocchi in memoria secondaria, ma semplifichiamo
assumendo che siano uguali:
\[
  1\; pagina \approx 1\; blocco
\] 

\subsubsection{Buffer}
Il buffer è una zona di memoria \textbf{condivisa dalle applicazioni} ed è importante
gestirlo in modo ottimizzato per massimizzare le prestazioni. Quando un dato viene
utilizzato più volte in tempi ravvicinati, il buffer evita di accedere alla memoria
secondaria. Il gestore del buffer si occupa di caricare dati dalla memoria secondaria
alla memoria centrale e viceversa ottimizzando gli accessi.

\subsubsection{Politica di gestione della memoria}
Le politiche della gestione della memoria sono:
\begin{itemize}
  \item \textbf{Lettura di un blocco}:
    \begin{itemize}
      \item Se il blocco \textbf{è già presente in una
        pagina del buffer}, allora non si esegue una lettura su memoria secondaria, ma
        si restituisce un puntatore alla pagina del buffer.

      \item Se il blocco \textbf{non è già presente in una pagina del buffer},
        si cerca una pagina libera e si carica il blocco nella pagina, restituendo
        il puntatore alla pagina stessa.
    \end{itemize}

  \item \textbf{Scrittura di un blocco}: In caso di richiesta di scrittura di un
    blocco precedentemente caricato in una pagina del buffer, il gestore del buffer
    può decidere di \textbf{differire la scrittura} su memoria secondaria in un secondo
    momento.
\end{itemize}
Queste politiche si basano sul \textbf{principio di località}, cioè i dati
referenziati di recente hanno maggiore probabilità di essere referenziati nuovamente
nel futuro.

Inoltre una nota \textbf{legge empirica} afferma che il 20\% dei dati viene referenziato
l'80\% delle volte.

\subsubsection{Gestione delle pagine}
Le pagine del buffer possono essere rappresentate tramite una tabella in cui
\begin{itemize}
  \item Per ogni \textbf{pagina del buffer} si memorizza il blocco \( B \) contenuto
    indicando il file e il numero di blocco (o offset). \( L \) indica che la pagina è
    libera

  \item Per ogni pagina del buffer si memorizza un insieme di \textbf{variabili di
    stato} tra cui si trovano sicuramente:
    \begin{itemize}
      \item Un \textbf{Contatore I} (\( B_I \)): indica il 
        \textbf{numero di transazioni} che stanno usando quella pagina

      \item Un \textbf{Bit di stato J} (\( B_J \)): indica se la pagina è stata
        modificata (1) o no (0)
    \end{itemize}
\end{itemize}
\begin{figure}[H]
  \centering
  \includegraphics[width=1\textwidth]{gestione_pagine}
  \caption{Gestione delle pagine del buffer}
\end{figure}

\subsubsection{Primitive per la gestione del buffer}
Le primitive per la gestione del buffer sono:
\begin{itemize}
  \item \texttt{fix}: viene usata dalle transazioni per richiedere l'accesso ad un
    blocco e restituisce un puntatore alla pagina contenente il blocco richiesto.
    I passi di questa primitiva sono:
    \begin{enumerate}
      \item Il blocco richiesto viene cercato nelle pagine del buffer, in caso sia
        presente, si restituisce un puntatore alla pagina contenente il blocco

      \item Altrimenti, viene scelta una pagina libera \( P \) (pagina \( L \) oppure
        con contatore \( I \) = zero). Si sceglie la pagina in base a criteri diversi:
        \begin{itemize}
          \item \textbf{LRU} (Least Recently Used): La pagina usata meno di recente
          \item \textbf{FIFO} (First In First Out): La pagina caricata da più tempo
        \end{itemize}
        \begin{itemize}
          \item 
            Se il bit di stato \( J \) di \( P \) è a 1 (cioè la pagina è stata modificata),
            \( P \) viene salvata in memoria secondaria (primitiva flush).

          \item Se \( J = 0 \) non si fa nulla
        \end{itemize}
            Si carica il blocco in P aggiornando file e numero di blocco corrispondenti.
            
      \item Se non esistono pagine libere, il gestore del buffer può adottare due
        politiche:
        \begin{itemize}
          \item \textbf{Steal}: "ruba" una pagina ad un’altra transazione
            (primitiva flush).
          \item \textbf{No steal}: sospende la transazione inserendola in una coda di
            attesa.
        \end{itemize}
        Se una pagina si libera (contatore \( I \) = zero) il gestore si comporterà
        come al punto precedente.

      \item Quando una transazione accede ad una pagina per la prima volta il contatore
        si incrementa.
    \end{enumerate}

  \item \texttt{unfix}: viene usata dalle transazioni per indicare che la transazione
    ha terminato di usare il blocco
    \begin{itemize}
      \item L'effetto è il decremento del contatore I che indica l'uso della pagina.
    \end{itemize}

  \item \texttt{setDirty}: viene usata dalle transazioni per indicare al gestore del
    buffer che il blocco della pagina è stato modificato
    \begin{itemize}
      \item L'effetto è la modifica del bit di stato \( J \) a 1
    \end{itemize}

  \item \texttt{flush}: viene usata dal gestore del buffer per salvare blocchi sulla
    bit di stato viene posto a zero)
    memoria secondaria in modo \textbf{asincrono}.
    \begin{itemize}
      \item L'effetto è la liberazione delle pagine "dirty", quindi il bit di stato
        \( J \) viene impostato a 0
    \end{itemize}

  \item \texttt{force}: viene usata per salvare in memoria secondaria in modo
    \textbf{sincrono} il blocco contenuto nella pagina (primitiva usata dal modulo
    Gestore dell’Affidabilità).
    \begin{itemize}
      \item L'effetto è il salvataggio in memoria secondaria del blocco e il bit di
        stato \( J \) posto a zero
    \end{itemize}
\end{itemize}

\subsection{Gestore dell'affidabilità}
Siccome il gestore del buffer differisce le operazioni di scrittura, questo potrebbe
essere pericoloso, quindi il gestore dell'affidabilità si occupa di garantire:
\begin{itemize}
  \item \textbf{Atomicità}: che le transazioni non vengano lasciate incomplete con
    alcune operazioni eseguite ed altre no

  \item \textbf{Persistenza}: che gli effetti di ciascuna transazione con conclusa
    con un commit siano mantenuti in modo permanente a prescindere dai guasti
    o dalle scritture differite
\end{itemize}

\subsubsection{File di log}
Il funzionamento si basa sull'utilizzo del \textbf{log}, cioè un archivio persistente
su cui vengono registrate tutte le operazioni eseguite dal DBMS e per il corretto
funzionamento il gestore deve disporre di un dispositivo di memoria \textbf{stabile},
cioè resistente ai guasti.
\begin{figure}[H]
  \centering
  \includegraphics[width=0.9\textwidth]{gestore_affidabilita_2}
  \caption{Gestore dell'affidabilità nell'architettura di un DBMS}
\end{figure}
\noindent
Il file di log viene memorizzato in memoria stabile e al suo interno vengono
memorizzati:
\begin{itemize}
  \item \textbf{Record di transazione}: cioè informazioni sulla transazione:
    \begin{itemize}
      \item \( B(T) \): Begin di una transazione \( T \) 
      \item \( C(T) \): Commit di una transazione \( T \)
      \item \( A(T) \): Abort di una transazione \( T \)
      \item Insert, delete, update eseguite da una transazione \( T \) su un oggetto
        \( O \), ad esempio una riga in una tabella
        \begin{itemize}
          \item \( I(T, O, AS) \): Insert di un oggetto \( O \) da parte di una
            transazione \( T \) con il suo \textbf{after state} \( AS \)
            (stato dell'oggetto dopo l'operazione)

          \item \( D(T, O, BS) \): Delete di un oggetto \( O \) da parte di una
            transazione \( T \) con il suo \textbf{before state} \( BS \)
            (stato dell'oggetto prima dell'operazione)

          \item \( U(T, O, BS, AS) \): Update di un oggetto \( O \) da parte di una
            transazione \( T \) con il suo \textbf{before state} \( BS \) e il suo
            \textbf{after state} \( AS \)
        \end{itemize}
    \end{itemize}
    L'obiettivo è quello di consentire il in caso di ripristino le seguenti operazioni:
    \begin{itemize}
      \item \textbf{UNDO}: per disfare un’azione su un oggetto \( O \)  è sufficiente
        ricopiare in \( O \) il valore \( BS \); l’insert/delete viene disfatto
        cancellando/inserendo \( O \);

      \item \textbf{REDO}: per rifare un’azione su un oggetto \( O \) è sufficiente
        ricopiare in \( O \) il valore \( AS \); l’insert/delete viene rifatto
        inserendo/cancellando \( O \)
    \end{itemize}
    Gli inserimenti controllano sempre l'esistenza dell'oggetto \( O \) prima di eseguire
    l'operazione, quindi non vengono inseriti i duplicati

    Per queste operazioni vale la \textbf{proprietà di idempotenza}, cioè:
    \[
      \begin{aligned}
        UNDO(UNDO(A)) &= UNDO(A) \\
        REDO(REDO(A)) &= REDO(A)
      \end{aligned}
    \] 

  \item \textbf{Record di sistema}: cioè informazioni sulle operazioni eseguite dal
    DBMS:
    \begin{itemize}
      \item \textbf{Operazione di DUMP} della base di dati: record di \texttt{DUMP}, copia
      completa e consistente della base di dati

      \item \textbf{Operazione di CheckPoint}: record \( CK(T_1, \ldots, T_n) \)
        indica che all'esecuzione del CheckPoint le transazioni attive erano
        \( T_1, \ldots, T_n \) 
        \begin{itemize}
          \item Sono elencate le transazioni che non hanno ancora espresso la loro
            scelta relativamente all'esito finale.

          \item Operazione svolta periodicamente dal gestore dell'affidabilità per
            garantire che gli effetti delle transazioni che hanno eseguito il commit
            siano registrati in modo permanente.

          \item Prevede i seguenti passi:
            \begin{enumerate}
              \item Sospensione delle operazioni di scrittura, commit e abort delle
                transazioni

              \item Esecuzione della primitiva \texttt{force} sulle pagine "dirty" di
                transazioni che hanno eseguito il commit

              \item Scrittura sincrona (primitiva \texttt{force}) sul file di LOG del
                record di CheckPoint con gli identificatori delle transazioni attive

              \item Ripresa delle operazioni di scrittura, commit e abort delle transazioni
            \end{enumerate}
        \end{itemize}
    \end{itemize}
\end{itemize}

\subsubsection{Regole di scrittura sul file di log}
Il gestore dell'affidabilità deve rispettare le seguenti regole di scrittura sul
file di log:
\begin{itemize}
  \item \textbf{Regola WAL} (Write-Ahead Logging): i record di log devono essere
    scritti sul LOG prima dell'esecuzione delle corrispondenti operazioni sulla base
    di dati (garantisce la possibilità di fare sempre UNDO)

  \item \textbf{Regola di Commit-Precedenza}: i record di log devono essere scritti
    sul LOG prima dell'esecuzione del COMMIT della transazione (garantisce la
    possibilità di fare sempre REDO)
\end{itemize}
Tali regole consentono inoltre di salvare i blocchi delle pagine "dirty" in modo
totalmente asincrono rispetto al commit delle transazioni.

\subsubsection{Guasti}
Una transazione \( T \) sceglie in modo atomico l'esito di COMMIT nel momento in
cui scrive nel file di log in modo sincrono (primitiva \texttt{force}) il suo
record di COMMIT \( C(T) \). Se avviene:
\begin{itemize}
  \item Un guasto prima della scrittura del record di COMMIT, allora si effettua un
    UNDO
  \item Un guasto dopo la scrittura del record di COMMIT, allora si effettua un REDO
\end{itemize}

\vspace{1em}
\noindent
Esistono due tipologie di guasto:
\begin{itemize}
  \item \textbf{Guasto di sistema}: è la perdita del contenuto della memoria centrale.
    Si risolve con la \textbf{ripresa a caldo} seguendo i seguenti passi:
    \begin{enumerate}
      \item Si accede al file di LOG e si ripercorre all'indietro il LOG fino al più
        recente record di CheckPoint

      \item Si inizializzano:
        \begin{itemize}
          \item L'insieme UNDO con tutte le transazioni da disfare
            \[
              UNDO = \text{tutte le transazioni presenti nel record di CheckPoint}
            \] 
          \item L'insieme REDO con tutte le transazioni da rifare
            \[
              REDO = \varnothing
            \] 
        \end{itemize}

      \item Si ripercorre in avanti il LOG e:
        \begin{itemize}
          \item Per ogni record \( B(T) \) incontrato si aggiunge \( T \) all'insieme
            UNDO

          \item Per ogni record \( C(T) \) incontrato si sposta \( T \) dall'insieme UNDO
            all'insieme REDO
        \end{itemize}

      \item Si ripercorre all'indietro il LOG disfacendo le operazioni eseguite dalle
        transazioni UNDO risalendo fino alla prima azione della transazione più vecchia.

      \item Si ripercorre in avanti il LOG per rifare le operazioni eseguite dalle
        transazioni REDO
    \end{enumerate}

  \item \textbf{Guasto di dispositivo}: è la perdita di parte o tutto il contenuto
    della base di dati in memoria secondaria. Si risolve con la ripresa a freddo
    seguendo i seguenti passi:
    \begin{enumerate}
      \item Si accede al DUMP della base di dati e si ricopia selettivamente la parte
        deteriorata della base di dati

      \item Si accede al LOG risalendo al record di DUMP

      \item Si ripercorre in avanti il LOG rieseguendo tutte le operazioni relative
        alla perte deteriorata comprese le azioni di commit e abort

      \item Si applica una ripresa a caldo
    \end{enumerate}
\end{itemize}

\begin{exercise}
  Data la seguente situazione del file di LOG:
  \[
    \begin{aligned}
      & B(T1), B(T2), U(T2,O1,B1,A1), I(T1,O2,A2), B(T3), C(T1), B(T4), \\
      &U(T3,O2,B3,A3), U(T4,O3,B4,A4), CK(T2,T3,T4), C(T4), B(T5), \\
      & U(T3,O3,B5,A5), U(T5,O4,B6,A6), D(T3,O5,B7), A(T3), C(T5), \\
      &I(T2,O6,A8) \text{ guasto}
    \end{aligned}
  \] 
  Che passi svolge la ripresa a caldo?

  \begin{enumerate}
    \item Risale il LOG fino all'ultimo record di CheckPoint \( CK(T2,T3,T4) \)

    \item Inizializza UNDO e REDO:
      \begin{itemize}
        \item \( UNDO = \left\{ T_2, T_3, T_4 \right\} \). Viene inizializzato con le
          transazioni del CheckPoint
        \item \( REDO = \varnothing \)
      \end{itemize}

    \item Ripercorre in avanti il LOG dal CheckPoint fino al guasto e si aggiungono
      o si spostano le transazioni tra UNDO e REDO:
      \begin{itemize}
        \item \( UNDO = \left\{ T_2, T_3, \cancel{T_4}, \cancel{\color{blue}T_5\color{black}} \right\}  \) 
        \item \( REDO = \left\{ T_4, T_5 \right\}  \) 
      \end{itemize}
      Quindi la situazione dopo questo passo è:
      \begin{itemize}
        \item \( UNDO = \left\{ T_2, T_3 \right\}  \)
        \item \( REDO = \left\{ T_4, T_5 \right\}  \)
      \end{itemize}

    \item Si ripercorre all'indietro il LOG disfacendo le operazioni eseguite dalle
      transazioni UNDO risalendo fino alla prima azione della transazione più vecchia.
      \[
        \begin{aligned}
          &DELETE(O6) \quad \text{(disfare insert)} \\
          &O5 := B7 \quad \text{(disfare delete)} \\
          &O3 := B5 \quad \text{(disfare update)} \\
          &O2 := B3 \quad \text{(disfare update)} \\
          &O1 := B1 \quad \text{(disfare update)}
        \end{aligned}
      \] 

    \item Si ripercorre in avanti il LOG per rifare le operazioni eseguite dalle
      transazioni REDO:
      \[
        \begin{aligned}
          &O3 := A4 \quad \text{(rifare update)} \\
          &O4 := A6 \quad \text{(rifare update)} \\
        \end{aligned}
      \]
  \end{enumerate}
\end{exercise}

\begin{exercise}
  Data la seguente situazione del file di LOG:
  \[
    \begin{aligned}
      &B(T1), B(T2), B(T3), I(T1,O1,A1), D(T2,O2,B2), B(T4), \\ 
      &U(T4,O3,B3,A3), U(T1,O4,B4,A4), C(T2), CK(T1,T3,T4), B(T5), B(T6), \\ 
      &U(T5,O5,B5,A5), A(T3), CK(T1,T4,T5,T6), B(T7), C(T4), \\
      &U(T7,O6,B6,A6), U(T6,O3,B7,A7), B(T8), A(T7), \text{ guasto}
    \end{aligned}
  \]
  Che passi svolge la ripresa a caldo?

  \begin{enumerate}
    \item Si risale il LOG fino all'ultimo record di CheckPoint \( CK(T1,T4,T5,T6) \)
    \item Inizializza UNDO e REDO:
      \begin{itemize}
        \item \( UNDO = \left\{ T1, T4, T5, T6 \right\} \)
        \item \( REDO = \varnothing \)
      \end{itemize}

    \item Ripercorre in avanti il LOG dal CheckPoint fino al guasto e si aggiungono
      o si spostano le transazioni tra UNDO e REDO:
      \begin{itemize}
        \item \( UNDO = \left\{ T1, \cancel{T4}, T5, T6, \color{blue}T7\color{black}, \color{blue}T8\color{black} \right\} \) 
        \item \( REDO = \left\{ T4 \right\} \) 
      \end{itemize}
      Quindi la situazione dopo questo passo è:
      \begin{itemize}
        \item \( UNDO = \left\{ T1, T5, T6, T7, T8 \right\} \) 
        \item \( REDO = \left\{ T4 \right\} \) 
      \end{itemize}

    \item Si ripercorre all'indietro il LOG disfacendo le operazioni eseguite dalle
      transazioni UNDO risalendo fino alla prima azione della transazione più vecchia:
      \[
        \begin{aligned}
          &O3 := B7 \quad \text{(disfare update)}\\
          &O6 := B6 \quad \text{(disfare update)}\\
          &O5 := B5 \quad \text{(disfare update)}\\
          &O4 := B4 \quad \text{(disfare update)}\\
          &DELETE(O1) \quad \text{(disfare insert)}
        \end{aligned}
      \] 

    \item Si ripercorre in avanti il LOG per rifare le operazioni eseguite dalle
      transazioni REDO:
      \[
        O3 := A3 \quad \text{(rifare update)}
      \]
  \end{enumerate}
  Si nota che ci sono due operazioni che riguardano l'oggetto \( O3 \) e l'ultima
  operazione è quella di \( T4 \) che ha eseguito il commit, quindi è quella che viene
  rifatta. Quella di \( T6 \) non viene applicata.
\end{exercise}

\subsection{Gestore dei metodi di accesso}
Il gestore dei metodi di accesso è il modulo che esegue il piano di esecuzione
prodotto dall'ottimizzatore e produce sequenze di richieste di accesso ai
blocchi della base di dati presenti in memoria secondaria. Le richieste
vengono inviate al gestore del buffer che si occupa di caricare i blocchi
necessari in pagine di memoria centrale.
\begin{figure}[H]
  \centering
  \includegraphics[width=0.6\textwidth]{richieste_metodi_accesso}
  \caption{Gestore dei metodi di accesso nell'architettura di un DBMS}
\end{figure}

\subsubsection{Metodi di accesso}
I metodi di accesso sono moduli \textbf{software} che implementano gli algoritmi di
\textbf{accesso e manipolazione dei dati} organizzati in specifiche strutture fisiche.
\begin{example}
  Alcuni esempi di metodi di accesso sono:
  \begin{itemize}
    \item Scansione sequenziale
    \item Accesso via indice
    \item Ordinamento
    \item Varie implementazioni del join
  \end{itemize}
\end{example}
Ogni metodo di accesso ai dati conosce:
\begin{itemize}
  \item \textbf{L'organizzazione delle tuple} (o dei record di indice) nei blocchi
    DATI (o INDICE) salvati in memoria secondaria. Come una tabella (o indice) viene
    organizzata in pagine di memoria secondaria.

  \item \textbf{L'organizzazione fisica interna delle pagine} sia quando contengono
    DATI, sia quando contengono strutture fisiche di accesso (o INDICI).
\end{itemize}

\subsubsection{Organizzazione di una pagina}
In una pagina sono presenti due tipi di informazioni:
\begin{itemize}
  \item \textbf{Informazioni utili}: dati veri e propri, cioè le tuple di una tabella

  \item \textbf{Informazioni di controllo}: consentono di accedere alle informazioni
    utili, ad esempio il dizionario, bit di parità, altre informazioni del file system,
    ecc.
\end{itemize}
La struttura di una pagina è la seguente:
\begin{figure}[H]
  \centering
  \includegraphics[width=0.9\textwidth]{organizzazione_pagina}
  \caption{Organizzazione di una pagina}
\end{figure}

\noindent
e i diversi campi sono:
\begin{itemize}
  \item \textbf{\textcolor{RedViolet}{Block header e block trailer}}:
    Ogni pagina, siccome coincide con un
    blocco di memoria di massa, ha una parte iniziale ed una parte finale che
    contengono informazioni di controllo utilizzate dal file system
    \begin{figure}[H]
      \centering
      \includegraphics[width=0.7\textwidth]{block_header_trailer}
      \caption{Block header e block trailer}
    \end{figure}

  \item \textbf{\textcolor{RedViolet}{Page header e page trailer}}:
    Ogni pagina, siccome contiene dati gestiti dal DBMS, ha una parte iniziale ed una
    finale contenenti l'informazione di controllo relativa alla specifica struttura
    fisica.
    Ad esempio:
    identificatore dell’oggetto (tabella, indice, dizionario dei dati, ecc.)
    contenuto nella pagina, puntatori a pagine successive o precedenti nella struttura
    dati, numero di dati utili elementari (tuple) contenuti nella pagina e quantità di
    memoria libera (contigua o non contigua) disponibile nella pagina.
    \begin{figure}[H]
      \centering
      \includegraphics[width=0.7\textwidth]{page_header_trailer}
      \caption{Page header e page trailer}
    \end{figure}

  \item 
    \begin{itemize}
      \item \textbf{\textcolor{RedViolet}{Dizionario di pagina}}: Contiene puntatori a
        ciascun dato utile elementare contenuto della pagina. Se le tuple sono a:
        \begin{itemize}
          \item \textbf{Lunghezza fissa}, allora il dizionario di pagina non è
            necessario perchè si deve soltanto memorizzare la dimensione delle tuple
            e l'offset dal punto iniziale
          \item \textbf{Lunghezza variabile} (valori nulli o stringhe), allora il
            dizionario memorizza l'offset di ogni tupla presente nel blocco e di
            ogni attributo di ogni tupla
        \end{itemize}
        Molti gestori non consentono la separazione di una tupla su più pagine.

      \item \textbf{\textcolor{RoyalBlue}{Parte utile}}: Contiene i dati veri e propri.
        Il dizionario di pagina e la parte utile crescono come stack contrapposti,
        lasciando libera la parte centrale in uno spazio contiguo.

      \item \textbf{\textcolor{Green}{Bit di parità}}: Verifica che l'informazione 
        contenuta nella pagina sia valida
    \end{itemize}
    \begin{figure}[H]
      \centering
      \includegraphics[width=0.7\textwidth]{dati_utili}
      \caption{Dizionario di pagina, parte utile e bit di parità}
    \end{figure}
\end{itemize}

\subsubsection{Operazioni sulle pagine}
Le operazioni che si possono eseguire sulle pagine sono:
\begin{itemize}
  \item \textbf{Inserimento ed aggiornamento di una tupla}:
    \begin{itemize}
      \item Se esiste spazio contiguo sufficiente: inserimento semplice.

      \item Se non esiste spazio contiguo ma esiste spazio sufficiente: riorganizzare lo
        spazio ed eseguire un inserimento semplice (operazione in memoria centrale).
        Operazione aggiuntiva: \textbf{riorganizzazione} della pagina.

      \item Se non esiste spazio sufficiente: interazione col file system per allocare
        nuovi blocchi per il file.
    \end{itemize}

  \item \textbf{Cancellazione}: è sempre possibile anche senza riorganizzare il
    contenuto nella pagina.

  \item \textbf{Accesso ad una tupla} particolare tramite il valore di chiave oppure
    l’offset nel dizionario.
    \begin{itemize}
      \item Accesso sequenziale a più tuple.
    \end{itemize}

  \item \textbf{Accesso ad un attributo di una tupla}: identificato in base all’offset
    e alla lunghezza del campo dopo aver identificato la tupla tramite la chiave o il
    suo offset.
\end{itemize}

\subsubsection{Strutture primarie per l'organizzazione di file}
La struttura primaria di un file stabilisce il criterio secondo il quale sono disposte
le tuple all'interno del file presente in memoria secondaria. Le strutture primarie più
comuni sono:
\begin{itemize}
  \item \textbf{Sequenziali}: disposizione consecutiva delle tuple che si basa su un
    criterio specifico, come l'ordine di inserimento

  \item \textbf{Ad accesso calcolato} (hash): collocano le tuple in posizioni
    determinate sulla base dell'esecuzione di un algoritmo di hashing

  \item \textbf{Ad albero}: utilizzate soprattutto come strutture secondarie (indici)
\end{itemize}

\begin{figure}[H]
  \centering
  \includegraphics[width=0.9\textwidth]{strutture_primarie}
  \caption{Strutture primarie per l'organizzazione di file}
\end{figure}

\subsubsection{Struttura sequenziale}
Un file è costituito da blocchi \textbf{logicamente} consecutivi e le tuple vengono
inserite nei blocchi rispettando una sequenza:
\begin{itemize}
  \item \textbf{Struttura sequenziale disordinata} (o seriale): le tuple sono disposte
    in base al loro ordine di inserimento, senza alcun criterio specifico. Questa è
    la struttura più semplice e più diffusa.

    L'efficienza di ogni operazione è la seguente:
    \begin{itemize}
      \item \textbf{Inserimento}: operazione molto efficiente, è sufficiente mantenere
        un riferimento all'ultimo blocco e richiede un solo accesso in memoria
        secondaria. Si avrà un costo maggiore in caso di verifica dei vincoli di
        integrità.

      \item \textbf{Ricerca}: operazione non efficiente, richiede ogni volta una
        scansione sequenziale che consideri tutti i record, con un costo lineare
        nel numero di blocchi. La ricerca si può fare usando strutture secondarie
        (indici).

      \item \textbf{Eliminazione}: facile una volta individuate le tuple
        coinvolte. Si marcano le tuple come "cancellate" senza riorganizzazione locale.

      \item \textbf{Modifica}: facile una volta individuate le tuple coinvolte. Si
        procede in loco se possibile, cioè quando non c'è stata nessuna modifica alla
        dimensione della tupla, oppure si procede con una cancellazione e un
        inserimento in fondo al file.
    \end{itemize}
    Eliminazioni e modifiche possono portare ad un uso non ottimale dello spazio
    siccome si devono effettuare riorganizzazioni periodiche.

  \item \textbf{Struttura sequenziale ad array}: le tuple sono disposte come in un
    array, e la loro posizione dipende dal valore assunto in ciascuna tupla da un
    campo detto \textbf{indice}. Questa struttura è utilizzabile soltanto quando
    le tuple di una tabella hanno una dimensione fissa, infatti non è quasi masi usata
    nei DBMS reali perchè difficilmente sono soddisfatte le condizioni per la sua
    applicazione.

    Ad ogni file viene assegnato un numero \( n \) di blocchi contigui e ciascun
    blocco è dotato di un numero \( m \) di posizioni disponibili per le tuple
    (quindi è un array bidimensionale \( n \times m \)). Ogni tupla è dotata di un
    valore numerico \( i \) che rappresenta l'indice della tupla in \( i \)-esima
    posizione nell'array.

    \begin{itemize}
      \item \textbf{Inserimento}: 
        \begin{itemize}
          \item Iniziale: l'indice \( i \) viene ottenuto incrementando un contatore
          \item Successivo: in una posizione libera oppure alla fine del file
        \end{itemize}

      \item \textbf{Cancellazioni}: creano delle posizioni libere
    \end{itemize}

  \item \textbf{Struttura sequenziale ordinata}: è un file sequenziale dove le tuple
    sono ordinate secondo una \textbf{chiave di ordinamento} (diversa dalla chiave
    primaria). Questa struttura rende efficienti le operazioni che hanno bisogno
    dell'ordinamento utilizzato, ad esempio:
    \begin{itemize}
      \item Elenco ordinato
      \item Range query
      \item Operazioni aggregate
    \end{itemize}
    In caso di aggiornamento è necessario mantenere l'ordinamento, quindi servono
    periodiche riorganizzazioni.

    \begin{itemize}
      \item \label{item:inserimento_ordinato} \textbf{Inserimento}: I passi da eseguire sono:
        \begin{enumerate}
          \item Individuare il blocco \( B \) che contiene \textbf{la tupla che precede},
            nell'ordine della chiave, la tupla da inserire
          \item Se \( B \) contiene spazio sufficiente per la nuova tupla allora si
            inserisce la tupla in \( B \)
          \item Altrimenti si aggiunge un nuovo blocco, detto \textbf{overflow page},
            alla struttura e si inserisce la tupla nel nuovo blocco e si aggiusta la
            catena di puntatori.
        \end{enumerate}

      \item \textbf{Scansione sequenziale ordinata secondo la chiave}: è sufficiente
        seguire i puntatori

      \item \label{item:cancellazione_ordinata} \textbf{Cancellazione di una tupla}: I passi da seguire sono:
        \begin{enumerate}
          \item Individuare il blocco \( B \) che contiene la tupla da cancellare
          \item Cancellare la tupla da \( B \)
          \item Aggiustare la catena di puntatori
        \end{enumerate}

      \item \textbf{Riorganizzazione}: si assegnano le tuple ai blocchi in base ad
        opportuni coefficienti di riempimento, riaggiustando i puntatori.
    \end{itemize}

    \begin{example}
      Il seguente esempio mostra una struttura sequenziale ordinata secondo la chiave
      \texttt{Filiale}:
      \begin{figure}[H]
        \centering
        \includegraphics[width=0.9\textwidth]{struttura_sequenziale_ordinata}
        \caption{Esempio di struttura sequenziale ordinata secondo la chiave \texttt{Filiale}}
      \end{figure}
    \end{example}
\end{itemize}

\subsubsection{Indici}
Per aumentare le prestazioni degli accessi alle tuple memorizzate nelle strutture
fisiche (\textbf{file sequenziale}), si introducono strutture ausiliarie dette
strutture di accesso ai dati, o \textbf{indici}. Queste strutture velocizzano l'accesso
casuale via chiave di ricerca, cioè un insieme di attributi utilizzati nell'indice
della ricerca. Ci sono due tipi di indici su file sequenziali:
\begin{itemize}
  \item \textbf{Indice primario}: La chiave di ordinamento del file sequenziale 
    \textbf{coincide} con la chiave di ricerca dell'indice.

    Ogni record dell'indice primario contiene una coppia \( \left<v_i, p_i \right> \),
    dove:
    \begin{itemize}
      \item \( v_i \) è il valore della chiave di ricerca
      \item \( p_i \) è il puntatore al primo record nel file sequenziale con chiave
        \( v_i \) 
    \end{itemize}
    Esistono due varianti dell'indice primario:
    \begin{itemize}
      \item \textbf{Indice denso}: per ogni occorrenza della chiave presente nel file
        esiste un corrispondente record nell'indice

      \item \textbf{Indice sparso}: solo per alcune occorrenze della chiave presenti
        nel file esiste un corrispondente record nell'indice, tipicamente una per
        blocco
    \end{itemize}

    Le operazioni che si possono eseguire su un indice primario sono:
    \begin{itemize}
      \item \textbf{Ricerca} di una tupla con chiave di ricerca \( K \).
        \begin{itemize}
          \item Se l'indice è \textbf{denso} \( K \) è presente nell'indice.
            \begin{enumerate}
              \item Scansione sequenziale dell'indice per trovare il record \( \left(K, p_k\right) \)
              \item Accesso al file attraverso il puntatore \( p_k \)
            \end{enumerate}
            Il costo è 1 accesso indice + 1 accesso blocco dati

            \begin{example}
              Un esempio di ricerca con indice primario denso è il seguente:
              (ricerca dei conti della filiale B)
              \begin{figure}[H]
                \centering
                \includegraphics[width=0.9\textwidth]{ricerca_indice_primario_denso}
                \caption{Esempio di ricerca con indice primario denso}
              \end{figure}
            \end{example}

          \item Se l'indice è \textbf{sparso} \( K \) potrebbe non essere presente
            nell'indice
            \begin{enumerate}
              \item Scansione sequenziale dell'indice fino al record
                \( \left( K', p_k \right) \), dove \( K' \) è il valore più grande che
                sia minore o uguale a \( K \)

              \item Accesso al file attraverso il puntatore \( p_k \) e scansione
                del file (blocco corrente) per trovare le tuple con chiave \( K \)
            \end{enumerate}
            Il costo è 1 accesso indice + 1 accesso blocco dati

            \begin{example}
              Un esempio di ricerca con indice primario sparso è il seguente:
              (ricerca dei conti della filiale B)
              \begin{figure}[H]
                \centering
                \includegraphics[width=0.9\textwidth]{ricerca_indice_primario_sparso}
                \caption{Esempio di ricerca con indice primario sparso}
              \end{figure}
            \end{example}


          \item \textbf{Inserimento} di un record nell'indice. Si effettuano gli
            stessi passi dell'inserimento in un file sequenziale
            (\ref{item:inserimento_ordinato}), solo che nel blocco di memoria secondaria
            al posto di tuple ci sono record di indice.

            \begin{itemize}
              \item Se l'indice è \textbf{denso} l'inserimento nell'indice avviene
                solo se la tupla inserita nel file ha un valore di chiave \( K \) che
                non è già presente

              \item Se l'indice è \textbf{sparso} l'inserimento avviene solo quando,
                per effetto dell'inserimento di una nuova tupla, si aggiunge un blocco
                dati alla struttura; in tutti gli altri casi l'indice rimane invariato
            \end{itemize}
            Gli indici sono costosi, quindi bisogna fare un compromesso sul numero di indici
            da mantenere e il numero di attributi ad accesso sequenziale.

          \item \textbf{Cancellazione} di un record nell'indice. Si effettuano gli stessi
            passi della cancellazione in un file sequenziale
            (\ref{item:cancellazione_ordinata}).

            \begin{itemize}
              \item Se l'indice è \textbf{denso} la cancellazione nell'indice avviene
                solo se la tupla cancellata nel file è l’ultima tupla con valore di
                chiave \( K \)  

              \item Se l'indice è \textbf{sparso} la cancellazione nell'indice avviene
                solo quando \( K \) è presente nell'indice e il corrispondente blocco
                viene eliminato; altrimenti, se il blocco sopravvive, va sostituito
                \( K \) nel record dell'indice con il primo valore \( K' \) presente
                nel blocco.
            \end{itemize}
        \end{itemize}
    \end{itemize}

  \item \textbf{Indice secondario}: La chiave di ordinamento e la chiave di ricerca
    sono diverse.

    Ogni record dell'indice secondario contiene una coppia \( \left<v_i, p_i \right> \),
    dove:
    \begin{itemize}
      \item \( v_i \) è il valore della chiave di ricerca
      \item \( p_i \) è il puntatore al \textbf{bucket di puntatori} che individuano
        nel file sequenziale tutte le tuple con chiave \( v_i \)
    \end{itemize}
    Gli indici secondari \textbf{sono sempre densi} perchè non si può sfruttare
    l'ordinamento delle tuple.

    Le operazioni che si possono eseguire su un indice secondario sono:
    \begin{itemize}
      \item \textbf{Ricerca} di una tupla con chiave di ricerca \( K \).
        \begin{enumerate}
          \item Scansione sequenziale dell'indice per trovare il record
            \( \left(K, p_k\right) \)
          \item Accesso al bucket \( B \) di puntatori attraverso il puntatore \( p_k \)
          \item Accesso al file attraverso i puntatori del bucket \( B \) 
        \end{enumerate}
        Il costo è 1 accesso indice + 1 accesso al bucket + \( n \) accessi alle pagine
        dati dove \( n \) è il numero di puntatori nel bucket \( B \)

        \begin{example}
          Un esempio di ricerca con indice secondario è il seguente:
          (ricerca dei conti di "Verdi")
          \begin{figure}[H]
            \centering
            \includegraphics[width=0.9\textwidth]{ricerca_indice_secondario}
            \caption{Esempio di ricerca con indice secondario}
          \end{figure}
        \end{example}

      \item \textbf{Inserimento e cancellazione}: si effettuano gli stessi passi
        dell'inserimento e della cancellazione in un indice primario denso, con in più
        l'aggiornamento dei bucket.
    \end{itemize}
\end{itemize}

\subsubsection{Struttura ad albero}
Quando l'indice aumenta di dimensioni, non si può più gestire in memoria centrale,
di conseguenza deve essere gestito in memoria secondaria. Per migliorare le prestazioni
degli accessi ad un indice di grandi dimensioni, si utilizza il B+-Tree, una struttura
ad albero con le seguenti caratteristiche:
\begin{itemize}
  \item Ogni \textbf{nodo viene memorizzato in una pagina/blocco di memoria secondaria}.

  \item I legami tra i nodi sono \textbf{puntatori a pagine di memoria secondaria}.

  \item Sono presenti \textbf{molti figli} e di conseguenza ha \textbf{pochi livelli}.

  \item L'albero è \textbf{bilanciato}, cioè la lunghezza dei percorsi che collegano
    la radice ai nodi foglia è costante

  \item Inserimenti e cancellazioni non alterano le prestazioni dell’accesso ai dati:
    l’albero si mantiene bilanciato
\end{itemize}

\subsubsection*{Struttura}
La struttura di un B+-Tree con \textbf{fan-out} \( = n \), cioè con \( n \) figli per
ogni nodo, è la seguente:
\begin{itemize}
  \item \textbf{Nodo foglia}:
    \begin{figure}[H]
      \centering
      \includegraphics[width=0.9\textwidth]{b_plus_tree}
      \caption{Struttura di un nodo foglia di un B+-Tree con fan-out \( = n \)}
    \end{figure}
    \noindent
    Può contenere fino a \( n-1 \) valori ordinati di chiave di ricerca e fino a \( n \) 
    puntatori ai dati. Ogni nodo foglia ha un \textbf{vincolo di ordinamento}, cioè
    \begin{itemize}
      \item Le chiavi dentro ad un nodo foglia sono ordinate e 
      \item I nodi foglia sono ordinati tra di loro. Dati due nodi foglia \( L_i \) e
        \( L_j \) con \( i < j \), allora:
        \[
          \forall k_t \in L_i; \forall k_s \in L_j; k_t < k_s
        \] 
    \end{itemize}
    Il puntatore \( p_m \) del nodo \( L_i \) punta al nodo foglia \( L_{i+1} \) se esiste:
    \begin{figure}[H]
      \centering
      \includegraphics[width=0.9\textwidth]{puntatore_foglia}
      \caption{Puntatore \( p_m \) del nodo \( L_i \) che punta al nodo \( L_{i+1} \)}
    \end{figure}

    Una variante del B+-Tree è quella in cui i nodi foglia contengono direttamente le
    tuple. In questo caso si ha una struttura fisica integrata dati/indice.

  \item \textbf{Nodo intermedio}: È una sequenza di \( m' \le n \) valori ordinati di
    chiave (può contenere fino a \( n \) puntatori a nodo). Ogni chiave \( K_i \) 
    è seguita da un puntatore \( p_i \):
    \begin{figure}[H]
      \centering
      \includegraphics[width=0.9\textwidth]{b_plus_tree_nodo_intermedio}
      \caption{Struttura di un nodo intermedio di un B+-Tree con fan-out \( = n \)}
    \end{figure}
\end{itemize}
Il fan-out è un parametro che fornisce un limite massimo, ma bisogna anche fornire
un limite minimo per evitare che i nodi siano troppo vuoti, questi limiti sono detti
\textbf{vincoli di riempimento}:
\begin{itemize}
  \item \textbf{Vincoli per i nodi foglia}: Ogni nodo foglia può contenere un numero
    di valori di chiave (\#chiavi) vincolato come:
    \[
      \left\lceil \frac{n-1}{2} \right\rceil \le \#chiavi \le n-1
    \] 

  \item \textbf{Vincoli per i nodi intermedi}: Ogni nodo intermedio può contenere
    un numero di puntatori (\#puntatori) vincolato come:
    \[
      \left\lceil \frac{n}{2} \right\rceil \le \#puntatori \le n
    \] 
\end{itemize}

\begin{example}
  Consideriamo ad esempio un B+-Tree con fan-out = 4 con contenuto:
  \[
    A, B, D, E, F, G, L, M, N, P
  \] 
  Il fan-out ci dice che per:
  \begin{itemize}
    \item I nodi foglia i vincoli di riempimento sono:
      \[
        2 \le \#chiavi \le 3
      \] 

    \item I nodi intermedi i vincoli di riempimento sono:
      \[
        2 \le \#puntatori \le 4
      \]

    \item La radice è l'unico nodo per cui i vincoli di riempimento \textbf{minimi}
      non si applicano
  \end{itemize}
  \begin{figure}[H]
    \centering
    \includegraphics[width=1\textwidth]{b_plus_tree_esempio}
    \caption{Esempio di B+-Tree con vincolo di riempimento minimo}
  \end{figure}
  \noindent
  La radice \( L \) punta:
  \begin{itemize}
    \item A sinistra: al sottoalbero con chiave \( K < "L" \) 
    \item A destra: al sottoalbero con chiave \( K \ge "L" \)
  \end{itemize}
  I figli puntano alle tuple con chiave uguale a quella del nodo foglia
  
  In questo caso l'albero per chiavi desiderate rispetta il vincolo di riempimento
  minimo. Per rispettare il vincolo di riempimento massimo bisogna fare un altro albero:
  \begin{figure}[H]
    \centering
    \includegraphics[width=1\textwidth]{b_plus_tree_esempio_2}
    \caption{Esempio di B+-Tree con vincolo di riempimento massimo}
  \end{figure}
  Le etichette per la radice sono assegnate secondo:
  \begin{itemize}
    \item Sinistra: Chiavi comprese tra: \( "E" \le K < "L" \) 
    \item Centro: Chiavi comprese tra \( "L" \le K < "N" \)
    \item Destra: Chiavi \( K \ge "N" \)
  \end{itemize}
\end{example}

\subsubsection*{Ricerca con chiave K}
Per effettuare la ricerca di una tupla con chiave di ricerca \( K \) si effettuano
i seguenti passi:
\begin{enumerate}
  \item \label{item:ricerca_b_plus_tree_1}
    Cercare nel \textbf{nodo radice} il più piccolo valore di chiave maggiore di
    \( K \)
    \begin{itemize}
      \item Se tale valore esiste (supponiamo sia \( K_i \)) allora seguire il
        puntatore \( p_i \) 
        \begin{figure}[H]
          \centering
          \includegraphics[width=0.8\textwidth]{ricerca_b_plus_tree}
          \caption{Ricerca di una tupla con chiave di ricerca \( K \) in un B+-Tree}
        \end{figure}

      \item Se tale valore non esiste seguire il puntatore \( p_m \)  
        \begin{figure}[H]
          \centering
          \includegraphics[width=0.8\textwidth]{ricerca_b_plus_tree_2}
          \caption{Ricerca di una tupla con chiave di ricerca \( K \) in un B+-Tree quando il più piccolo valore di chiave maggiore di \( K \) non esiste}
        \end{figure}
    \end{itemize}

  \item Se il nodo raggiunto è un \textbf{nodo foglia} cercare il valore \( K \) al suo
    interno e seguire il corrispondente puntatore verso le tuple, se non esistono
    riprendere il passo \ref{item:ricerca_b_plus_tree_1}
    \begin{figure}[H]
      \centering
      \includegraphics[width=0.8\textwidth]{ricerca_b_plus_tree_3}
      \caption{Ricerca di una tupla con chiave di ricerca \( K \) in un B+-Tree quando si raggiunge un nodo foglia}
    \end{figure}
\end{enumerate}

\begin{example}
  Ad esempio ricerchiamo la tupla con chiave di ricerca \( K = "G" \):
  \begin{figure}[H]
    \centering
    \includegraphics[width=1\textwidth]{ricerca_b_plus_tree_esempio}
    \caption{Esempio di ricerca di una tupla con chiave di ricerca \( K \) in un B+-Tree}
  \end{figure}
\end{example}

\vspace{1em}
\noindent
Il costo della ricerca dipende dalla profondità dell'albero in termini di
\textbf{accessi alla memoria secondaria}, dove ogni nodo rappresenta un blocco. 
Il costo in una struttura ad albero è pari alla profondità dell'albero, che nel
B+-Tree è indipendente dal percorso ed è funzione del fan-out \( n \) e del numero
di valori chiave presenti nell'albero \#valoriChiave
\[
  prof_{\text{B+-Tree}} \le \underbrace{1}_{\text{Radice}} + 
  \underbrace{\log_{\left\lceil \frac{n}{2} \right\rceil} \left( 
      \overbrace{\frac{\#valoriChiave}{\left\lceil \frac{n-1}{2} \right\rceil }}^{\text{Nodi foglia massimi}}
  \right)}_{\text{\# Livelli}}
\] 

\subsubsection*{Inserimento}
Per inserire una tupla in un B+-Tree bisogna mantenere i vincoli di riempimento, quindi
vengono eseguiti i seguenti passi:
\begin{enumerate}
  \item \textbf{Ricerca} del nodo foglia \( NF \) dove \( K \) va inserito
  \item 
  \begin{itemize}
    \item Se \( K \) è presente in \( NF \):
      \begin{itemize}
        \item Se è \textbf{indice primario}: non si fa nessuna azione
        \item Se è \textbf{indice secondario}: si aggiornano i bucket di puntatori
      \end{itemize}

    \item Altrimenti si inserisce \( K \) in \( NF \) rispettando l'ordine:
      \begin{itemize}
        \item Se è \textbf{indice primario}: si inserisce \( K \) e si aggiorna il
          puntatore alla prima tupla

        \item Se è \textbf{indice secondario}: si inserisce \( K \) e si crea
          il bucket
      \end{itemize}
      Se all'inserimento di \( K \) viene violato il vincolo di riempimento massimo,
      bisogna eseguire uno \textbf{split} del nodo \( NF \) in due nodi foglia.
      Lo split viene eseguito se in \( NF \) esistono già \( n \) valori di chiave
      e bisogna eseguire i seguenti passi:
      \begin{enumerate}
        \item Creare due nodi foglia \( NF_1 \) e \( NF_2 \) 

        \item Inserire i primi \( \left\lceil \frac{n-1}{2} \right\rceil  \) valori
          di chiave in \( NF_1 \)

        \item Inserire i restanti in \( NF_2 \)

        \item Aggiornare il padre, cioè se il padre viola il vincolo di riempimento
          massimo bisogna propagare lo split e quindi creare una nuova radice/livello
          se necessario
      \end{enumerate}
  \end{itemize}
\end{enumerate}

\begin{example}
  Un esempio di split di un nodo foglia con fan-out = 5 è il seguente:
  \begin{figure}[H]
    \centering
    \includegraphics[width=1\textwidth]{split_b_plus_tree}
    \caption{Esempio di split di un nodo foglia in un B+-Tree con fan-out = 5}
  \end{figure}
  \noindent
  Se si aggiungesse la chiave \( 15 \) tra il \( 12 \) e il \( 16 \)  si violerebbe
  il vincolo di riempimento massimo: \( \#chiavi = 5 \), quindi si esegue lo split
  del nodo foglia.
\end{example}

\subsubsection*{Cancellazione}
Per cancellare una tupla in un B+-Tree con chiavee \( K \) bisogna mantenere i vincoli
di riempimento, quindi vengono eseguiti i seguenti passi:
\begin{enumerate}
  \item \textbf{Ricercare} il nodo foglia \( NF \) che contiene \( K \)
  \item Cancellare \( K \) e il suo puntatore. Potrebbe violare il vincolo di riempimento minimo, in quel caso bisogna
        fare il \textbf{merge} di \( NF \) con un fratello.
        Per fare il merge, se nel nodo da unire esistono \( \left\lceil \frac{n-1}{2} \right\rceil -1 \) 
        valori chiave, si procede come segue:
        \begin{enumerate}
          \item Individuare il nodo fratello
          \item Si genera un unico nodo foglia
          \item Si aggiorna il padre
            \begin{itemize}
              \item Se non c'è un nodo foglia fratello che unito rispetta il vincolo
                di riempimento massimo, si fa prima il merge e poi lo split

              \item Se anche il padre viola il vincolo di riempimento massimo, si
                propaga il merge fino ad arrivare alla radice
            \end{itemize}
        \end{enumerate}
\end{enumerate}

\begin{example}
  Un esempio di merge di un nodo foglia con fan-out = 5, in cui si vuole eliminare
  la tupla con chiave \( K = "3" \) è il seguente:
  \begin{figure}[H]
    \centering
    \includegraphics[width=1\textwidth]{merge_b_plus_tree}
    \caption{Esempio di merge di un nodo foglia in un B+-Tree con fan-out = 5}
  \end{figure}
  \noindent
  Se si eliminasse il nodo con chiave \( 3 \) si violerebbe il vincolo di riempimento
  minimo, quindi si deve fare il merge con il nodo fretello.
\end{example}

\subsubsection{Struttura ad accesso calcolato (hash)}
Questa struttura garantisce un \textbf{accesso associativo} alle tuple, cioè la
locazione fisica dei dati dipende dal valore assunto dalla chiave. Si potrebbe
utilizzare un array indicizzato dalla chiave, ma questo non è pratico perchè i valori
della chiave sono tanti, quindi si utilizza una \textbf{funzione di hash} che mappa
i valori di chiave in indirizzi di memorizzazione delle tuple:
\[
  h: \; K \to B
\] 
dove:
\begin{itemize}
  \item \( K \) è il dominio delle chiavi
  \item \( B \) è il dominio degli indirizzi di memorizzazione
\end{itemize}
Il vantaggio di questa struttura è che il costo della ricerca è costante e uguale
ad un accesso alla memoria secondaria. Lo svantaggio è quello delle \textbf{collisioni},
cioè valori di chiavi diversi possono produrre lo stesso indirizzo. Una funzione
di hash \textbf{buona} minimizza il numero di collisioni.

\begin{definition}
  Il fattore di riempimento definisce la frazione di spazio fisico mediamente
  utilizzata da un blocco
\end{definition}

\noindent
Siano:
\begin{itemize}
  \item \( T \) il numero di tuple
  \item \( F \) il fattore blocco, cioè il numero di tuple che ci stanno in un blocco
  \item \( f \) il fattore di riempimento
\end{itemize}
Si ha che il numero di blocchi è definito come:
\[
  \#B = \left\lceil \frac{T}{F \cdot f} \right\rceil 
\] 
Il DBMS alloca un numero di bucket di puntatori uguale al numero di blocchi \( \#B \).
Si definiscono due funzioni:
\begin{itemize}
  \item Funzione di \textbf{folding}: trasforma i valori chiave in numeri interi
    positivi
    \[
      f: \; K \to \mathbb{Z}^+
    \] 

  \item Funzione di \textbf{hashing}:
    \[
      h: \; \mathbb{Z}^+ \to B
    \] 
\end{itemize}

\subsubsection*{Ricerca con chiave K}
Per effettuare la ricerca di una tupla con chiave di ricerca \( K \) si effettuano
i seguenti passi:
\begin{enumerate}
  \item Calcolare \( b = h(f(K)) \) (costo zero)
  \item Accedere al bucket \( b \) (costo: 1 accesso alla pagina di memoria secondaria)
  \item Accedere alle \( n \) tuple attraverso i puntatori del bucket (costo:
    \( m \) accessi a pagine di memoria secondaria, con \( m \le n \))
\end{enumerate}

\begin{example}
  Un esempio di creazione di una tabella di hash è il seguente:
  \begin{figure}[H]
    \centering
    \includegraphics[width=1\textwidth]{hash_table_esempio}
    \caption{Esempio di creazione di una tabella di hash}
  \end{figure}
\end{example}

% \section{Transazioni concorrenti}
% Il componente dell'architettura del DBMS che si occupa di gestire le transazioni
% concorrenti è il \textbf{gestore della concorrenza} che riceve richieste di accesso
% ai dati e decide se autorizzarle o meno, eventualmente riordinandole, come se
% fosse uno scheduler.
% \begin{figure}[H]
%   \centering
%   \includegraphics[width=0.9\textwidth]{gestore_concorrenza}
%   \caption{Gestore della concorrenza nell'architettura di un DBMS}
% \end{figure}
%
% \noindent
% L'unità di misura per caratterizzare il carico applicativo di un DBMS è il numero di
% transazioni al secondo (TPS, Transactions Per Second). Per gestire un carico elevato,
% tipico delle applicazioni gestionali (100 o 1000 tps), un DBMS deve eseguire le
% transazioni in modo concorrente.
%
% \subsection{Anomalie}
% L'esecuzione concorrente di transazioni senza controllo può generare \textbf{anomalie}
% o problemi di correttezza ed è quindi necessario introdurre dei meccanismi di
% controllo nell'esecuzione delle transazioni per evitare queste anomalie.
%
% \begin{define}[Notazione]
%   Verrà utilizzata la seguente notazione:
%   \begin{itemize}
%     \item \( t_i \) indica una transazione
%     \item \( r_i(x) \) è un operazione di lettura eseguita dalla transazione \( t_i \) 
%       sulla risorsa \( x \) 
%     \item \( w_i(x) \) è un operazione di scrittura eseguita dalla transazione \( t_i \) 
%       sulla risorsa \( x \)
%
%     \item \texttt{bot} indica l'inizio di una transazione (begin of transaction)
%     \item \texttt{commit} indica l'applicazione definitiva degli effetti di una
%       transazione (commit)
%     \item \texttt{eot} indica la fine di una transazione (end of transaction)
%   \end{itemize}
% \end{define}
%
% \subsubsection{Perdita di aggiornamento}
% Questa anomalia si verifica quando gli effetti di una transazione concorrente sono
% persi. Si considerino le seguenti transazioni:
% \[
% \begin{aligned}
%   t_1: & bot \; r_1(x); \; x = x + 1; \; w_1(x); \; commit; \; eot \\
%   t_2: & bot \; r_2(x); \; x = x + 1; \; w_2(x); \; commit; \; eot
% \end{aligned}
% \] 
% L'esecuzione concorrente di \( t_1 \) e \( t_2 \) può portare alla perdita di
% aggiornamento, come si può vedere dalla seguente simulazione:
% \begin{figure}[H]
%   \centering
%   \includegraphics[width=1\textwidth]{perdita_aggiornamento}
%   \caption{Esempio di perdita di aggiornamento}
% \end{figure}
%
% \subsubsection{Lettura inconsistente}
% Questa anomalia si verifica quando accessi successivi ad uno stesso dato, all'interno
% di una transazione, ritornano valori diversi. Si considerino le seguenti transazioni:
% \[
% \begin{aligned}
%   t_1: & bot \; r_1(x); \; r_1'(x); \; commit; \; eot \\
%   t_2: & bot \; r_2(x); \; x = x + 1; \; w_2(x); \; commit; \; eot
% \end{aligned}
% \] 
% L'esecuzione concorrente di \( t_1 \) e \( t_2 \) può portare alla lettura
% inconsistente, come si può vedere dalla seguente simulazione:
% \begin{figure}[H]
%   \centering
%   \includegraphics[width=1\textwidth]{lettura_inconsistente}
%   \caption{Esempio di lettura inconsistente}
% \end{figure}
%
% \subsubsection{Lettura sporca}
% Questa anomalia si verifica quando viene letto un dato che rappresenta uno stato
% intermedio nell’evoluzione di una transazione. Si considerino le seguenti transazioni:
% \[
% \begin{aligned}
%   t_1: & bot \; r_1(x); \; commit; \; eot \\
%   t_2: & bot \; r_2(x); \; x = x + 1; \; w_2(x); \; \ldots abort;
% \end{aligned}
% \] 
% L'esecuzione concorrente di \( t_1 \) e \( t_2 \) può portare alla lettura
% sporca, come si può vedere dalla seguente simulazione:
% \begin{figure}[H]
%   \centering
%   \includegraphics[width=1\textwidth]{lettura_sporca}
%   \caption{Esempio di lettura sporca}
% \end{figure}
%
% \subsubsection{Aggiornamento fantasma}
% Questa anomalia si verifica quando avviene l'osservazione di uno stato intermedio che
% non soddisfa i vincoli di integrità. Si considerino le seguenti transazioni:
% \[
% \begin{aligned}
%   Risorse &: x, y, z \\
%   Vincoli &: x + y + z = 100 \\
% \end{aligned}
% \] 
% \[
% \begin{aligned}
%   t_1: & bot \; r_1(y); \; r_1(x); \; r_1(z); \; s = x + y + z; \; eot \\
%   t_2: & bot \; r_2(y); \; y = y + 10; \; r_2(z); \; z = z - 10; \; w_2(y); \; w_2(z);
%   \; commit; \; eot
% \end{aligned}
% \] 
% L'esecuzione concorrente di \( t_1 \) e \( t_2 \) può portare all'aggiornamento
% fantasma, come si può vedere dalla seguente simulazione:
% \begin{figure}[H]
%   \centering
%   \includegraphics[width=1\textwidth]{aggiornamento_fantasma}
%   \caption{Esempio di aggiornamento fantasma}
% \end{figure}
%
% \subsubsection{Inserimento fantasma}
% Questa anomalia si verifica quando avvengono valutazioni diverse di valori aggregati a
% causa di inserimenti intermedi.
% La transazione che valuta un valore aggregato relativo a tutti gli elementi che
% soddisfano un predicato di selezione (ad esempio il voto medio degli studenti del terzo
% anno) può subire l'anomalia di inserimento fantasma se il valore aggregato viene
% valutato due volte e tra la prima e la seconda valutazione viene inserito un nuovo
% studente, quindi i due valori medi letti potrebbero essere diversi. Questa anomalia
% non si può risolvere facendo riferimento solo ai dati già presenti nella base di dati,
% ma è necessario tener presente che potrebbero esserci delle nuove tuple che compaiono
% \textit{improvvisamente}.
%
% \subsection{Schedule}
% Uno schedule è una sequenza di operazioni di lettura e scrittura eseguite su risorse
% della base di dati da diverse transazioni concorrenti che rappresenta una possibile
% esecuzione concorrente di diverse transazioni:
% \[
%   S: r_1(x); \; r_2(z); \; w_1(x); \; w_2(z);
% \] 
% Per evitare anomalie si stabilisce di accettare \textbf{solo} gli schedule equivalenti
% ad uno schedule seriale.
%
% \subsubsection{Schedule seriale}
% Lo schedule seriale è uno schedule dove le operazioni di ogni
% transazione compaiono in sequenza \textbf{senza essere inframezzate} da operazioni
% di altre transazioni.
% \begin{example}
%   Consideriamo i seguenti esempi di schedule:
%   \begin{itemize}
%     \item \( S_1: r_1(x); \; \color{blue}r_2(x); \; w_2(x); \; \color{black}w_1(x); \)
%       \textbf{non} è seriale
%
%     \item \( S_2: r_1(x); \; w_1(x); \; \color{blue} r_2(x); \; w_2(x); \) 
%       \textbf{è} seriale
%   \end{itemize}
% \end{example}
%
%
% \subsubsection{Schedule serializzabile}
% Lo schedule serializzabile è uno schedule equivalente ad uno schedule
% seriale. Due schedule sono equivalenti se \textbf{producono gli stessi effetti}
% sulla base di dati.
%
% \subsection{Tecniche di gestione della concorrenza}
% Per le seguenti tecniche di gestione della concorrenza, si assume l'ipotesi di
% \textbf{commit-proiezione}, cioè si suppone che le transazioni abbiano esito noto.
% Quindi si tolgono dagli schedule tutte le operazioni delle transazioni che non vanno
% a buon fine.
%
% \subsubsection{View-equivalenza}
% \begin{definition}[Relazione LEGGE\_DA]
%   Dato uno schedule \( S \) si dice che un'operazione di lettura \( r_i(x) \), che
%   compare in \( S \), \texttt{LEGGE\_DA} un'operazione di scrittura \( w_j(x) \), che
%   compare in \( S \), se \( w_j(x) \) \textbf{precede} \( r_i(x) \) in \( S \) e non
%   vi è alcuna operazione di scrittura \( w_k(x) \) tra le due.
%
%   \begin{example}
%     Un esempio di utilizzo della relazione LEGGE\_DA è il seguente:
%     \begin{itemize}
%       \item \( S_1: r_1(x);\; r_2(x);\; w_2(x);\; w_1(x);\, \) 
%         \texttt{LEGGE\_DA}\( (S_1) = \varnothing \) 
%
%       \item \( S_2: r_1(x);\; w_1(x);\; r_2(x);\; w_2(x);\, \) 
%         \texttt{LEGGE\_DA}\( (S_2) = \{ (r_2(x), w_1(x)) \} \)
%     \end{itemize}
%   \end{example}
% \end{definition}
%
% \begin{definition}[SCRITTURE\_FINALI]
%   Dato uno schedule \( S \) si dice che un'operazione di scrittura \( w_i(x) \), che
%   compare in \( S \), è una \textbf{SCRITTURA FINALE} se è l'ultima operazione di
%   scrittura della risorsa \( x \) in \( S \).
%
%   \begin{example}
%     Un esempio di utilizzo della relazione \\ SCRITTURE\_FINALI è il seguente:
%     \begin{itemize}
%       \item \( S_1: r_1(x);\; r_2(x);\; w_2(x);\; w_3(y);\; w_1(x); \) \\
%         \texttt{SCRITTURE\_FINALI}\( (S_1) = \{ w_3(y), w_1(x) \} \)
%
%       \item \( S_2: r_1(x);\; w_1(x);\; r_2(x);\; w_2(x);\; w_3(y); \) \\
%         \texttt{SCRITTURE\_FINALI}\( (S_2) = \{ w_2(x), w_3(y) \} \)
%     \end{itemize}
%   \end{example}
% \end{definition}
%
% \begin{definition}[View-equivalenza]
%   Due schedule \( S_1 \) e \( S_2 \) sono view-equivalenti (\( S_1 \approx_V S_2 \))
%   se possiedono le stesse relazioni LEGGE\_DA e SCRITTURE\_FINALI.
% \end{definition}
%
% \begin{definition}[View-serializzabilità]
%   Uno schedule \( S \) è view-serializzabile (VSR) se esiste uno schedule seriale \( S' \)
%   tale che \( S \approx_V S' \).
%
%   \vspace{1em}
%   \noindent
%   \textbf{Attenzione:} gli schedule seriali \( S' \) si generano considerando tutte le
%   possibili permutazioni delle transazioni che compaiono in \( S \). Inoltre \textbf{non
%   si deve} cambiare l'ordine delle operazioni di ciascuna transazione perchè ciò
%   significherebbe modificare la transazione.
% \end{definition}
%
% \noindent
% Gli schedule corrispondenti alle anomalie di:
% \begin{itemize}
%   \item Perdita di aggiornamento
%   \item Letture inconsistenti
%   \item Aggiornamenti fantasma
% \end{itemize}
% \textbf{non sono view-serializzabili}. Inoltre l'algoritmo per il test di view-equivalenza
% tra due schedule è di complessità lineare, mentre l'algoritmo per il test di
% view-serializzabilità di uno schedule è di complessità esponenziale.
%
% \vspace{1em}
% \noindent
% In conclusione si è osservato che la view-serializzabilità richiede l'ipotesi di
% commit-proiezione, e algoritmi di complessità troppo elevata, quindi non è applicabile
% nei sistemi reali.
%
% \subsubsection{Conflict-equivalenza}
%
% \begin{definition}[Conflitto]
%   Dato uno schedule \( S \) si dice che una coppia di operazioni \( (a_i, a_j) \),
%   dove \( a_i \) e \( a_j \) compaiono nello schedule \( S \), rappresentano un
%   conflitto se:
%   \begin{itemize}
%     \item appartengono a transazioni diverse (\( i \neq j \))
%
%     \item entrambe operano sulla stessa risorsa
%
%     \item almeno una di esse è un'operazione di scrittura
%
%     \item \( a_i \) compare in \( S \) prima di \( a_j \)
%   \end{itemize}
% \end{definition}
%
% \begin{definition}[Conflict-equivalenza]
%   Due schedule \( S_1 \) e \( S_2 \) sono conflict-equivalenti (\( S_1 \approx_C S_2 \))
%   se possiedono le stesse operazioni e gli stessi conflitti, cioè se le operazioni 
%   in conflitto sono nello stesso ordine nei due schedule.
% \end{definition}
%
% \begin{definition}[Conflict-serializzabilità]
%   Uno schedule \( S \) è \\ conflict-serializzabile (CSR) se esiste uno schedule seriale
%   \( S' \) tale che \( S \approx_C S' \).
% \end{definition}
%
% \noindent
% L'algoritmo per il test di conflict-serializzabilità di uno schedule è di \textbf{
% complessità lineare} ed è il seguente:
%
% Dato uno schedule \( S \):
% \begin{itemize}
%   \item Si costruisce il grafo dei conflitti \( G(N, A) \) dove:
%     \begin{itemize}
%       \item \( N = \left\{ t_1, \ldots, t_n \right\} \) è l'insieme delle transazioni
%         di \( S \) 
%
%       \item \( (t_i, t_j) \in A \) se esiste almeno un conflitto \( (a_i, a_j) \) in 
%         \( S \)
%     \end{itemize}
%
%   \item Se il grafo è \textbf{aciclico}, allora \( S \) è conflict-serializzabile.
% \end{itemize}
% Richiede comunque l'ipotesi di commit-proiezione.
%
% \begin{theorem}
%   La conflict-serializzabilità è una condizione sufficiente, ma non necessaria per
%   la view-serializzabilità, cioè \( CSR \subset VSR \):
%   \begin{figure}[H]
%     \centering
%     \includegraphics[width=0.6\textwidth]{csr_vsr}
%     \caption{Relazione tra conflict-serializzabilità e view-serializzabilità}
%   \end{figure}
% \end{theorem}
%
% \subsubsection{Locking a due fasi stretto}
% Il locking a due fasi (2PL, Two-Phase Locking) è il metodo applicato nei sistemi reali
% per la gestione dell'esecuzione concorrente delle transazioni perchè non richiede di
% conoscere in anticipo l'esito delle transazioni. Il locking a due fasi è un sottoinsieme
% della conflict-serializzabilità.
% Gli aspetti che caratterizzano il locking a due fasi sono i seguenti:
% \begin{itemize}
%   \item \textbf{Meccanismo di base per la gestione dei lock}:
%     Si basa sull'introduzione di alcune \textbf{primitive di lock} che consentono
%     alle transazioni di bloccare le risorse sulle quali vogliono agire con operazioni
%     di lettura e scrittura. Le primitive di lock sono le seguenti:
%     \begin{itemize}
%       \item \( r\_lock_K(x) \): richiesta di un lock \textbf{condiviso} da parte della
%         transazione \( t_K \) sulla risorsa \( x \) per eseguire una lettura
%
%       \item \( w\_lock_K(x) \): richiesta di un lock \textbf{esclusivo} da parte della
%         transazione \( t_K \) sulla risorsa \( x \) per eseguire una scrittura
%
%       \item \( unlock_K(x) \): richiesta da parte della transazione \( t_K \) di
%         liberare la risorsa \( x \) da un precedente lock
%     \end{itemize}
%     Le transazioni devono seguire le seguenti regole per applicare le primitive di
%     lock:
%     \begin{itemize}
%       \item Ogni \textbf{lettura} deve essere preceduta da un \( r\_lock \) e seguita
%         da un \( unlock \). Sono ammessi più \( r\_lock \) contemporanei sulla stessa
%         risorsa (lock condiviso)
%
%       \item Ogni \textbf{scrittura} deve essere preceduta da un \( w\_lock \) e seguita
%         da un \( unlock \). \textbf{Non} sono ammessi più \( w\_lock \) (oppure
%         \( w\_lock \) e \( r\_lock \)) contemporanei sulla stessa risorsa (lock
%         esclusivo)
%     \end{itemize}
%     Se una transazione segue le precedenti regole, allora si dice \textbf{ben formata}
%     rispetto al locking.
%
%   \item \textbf{Politica di concessione dei lock sulle risorse}:
%     Il gestore dei lock (o gestore della concorrenza) mantiene per ogni risorsa \( x \) 
%     le seguenti informazioni:
%     \begin{itemize}
%       \item \textbf{Stato}: \( s(x) \in \left\{ libero, r\_lock, w\_lock \right\} \) 
%
%       \item \textbf{Transazioni in r\_lock}: \( c(x) = \left( t_1, \ldots, t_n \right)  \),
%         dove \( t_1, \ldots, t_n \) hanno un \( r\_lock \) su \( x \)
%     \end{itemize}
%     Il comportamento del gestore dei lock è descritto dalla seguente tablella:
%     \begin{table}[H]
%       \centering
%       \includegraphics[width=0.9\textwidth]{gestore_lock}
%       \caption{Comportamento del gestore dei lock}
%     \end{table}
%
%   \item \textbf{Regola che garantisce la serializzabilità}:
%     Una transazione \textbf{dopo aver rilasciato un lock non può acquisirne altri}.
%     Questa regola è visualizzata dal seguente grafico:
%     \begin{figure}[H]
%       \centering
%       \includegraphics[width=0.9\textwidth]{grafico_locking_due_fasi}
%       \caption{Locking a due fasi}
%     \end{figure}
% \end{itemize}
%
% \noindent
% Il locking a due fasi stretto (Strict 2PL) è una variante del locking a due fasi che
% introduce una condizione aggiuntiva per rimuovere la necessità dell'ipotesi di
% commit-proiezione. La condizione aggiuntiva è la seguente: Una transazione può
% rilasciare i lock solo quando ha eseguito correttamente un \textbf{commit} o un
% \textbf{rollback}.
% \begin{figure}[H]
%   \centering
%   \includegraphics[width=0.9\textwidth]{grafico_locking_due_fasi_stretto}
%   \caption{Locking a due fasi stretto}
% \end{figure}
%
% \subsubsection{Blocco critico (deadlock)}
% Il blocco critico è una situazione di blocco che si verifica nell'esecuzione di
% transazioni concorrenti quando:
% \begin{itemize}
%   \item Due transazioni hanno bloccato delle risorse:
%     \begin{itemize}
%       \item \( t_1 \) ha bloccato \( r_1 \)
%
%       \item \( t_2 \) ha bloccato \( r_2 \)
%     \end{itemize}
%
%   \item Inoltre \( t_1 \) è in attesa sulla risorsa \( r_2 \) e \( t_2 \) è in attesa
%     sulla risorsa \( r_1 \)
% \end{itemize}
% Se il numero medio di tuple per tabella è \( n \) e la granularità del lock (cioè
% il numero di risorse che possono essere bloccate) è la tupla, la probabilità che
% si verifichi un lock tra due transazioni è:
% \[
%   P(\text{deadlock di lunghezza 2}) = \frac{1}{n^2}
% \] 
% \begin{example}
%   Un esempio di deadlock utilizzando il locking a due fasi stretto è il seguente:
%   \begin{figure}[H]
%     \centering
%     \includegraphics[width=0.9\textwidth]{deadlock_locking_due_fasi_stretto}
%     \caption{Esempio di deadlock con locking a due fasi stretto}
%   \end{figure}
% \end{example}
%
% \noindent
% Le tecniche per risolvere i deadlock sono le seguenti:
% \begin{itemize}
%   \item \textbf{Timeout}: se una transazione è in attesa su una risorsa da troppo tempo 
%     viene abortita
%
%   \item \textbf{Prevenzione}: 
%     \begin{itemize}
%       \item Una transazione blocca tutte le risorse a cui deve accedere in lettura o
%         scrittura in una sola volta (o blocca tutto o non blocca nulla)
%
%       \item Ogni transazione \( t_i \) acquisisce un timestamp \( TS_i \) all'inizio
%         della sua esecuzione. La transazione \( t_i \) può attendere una risorsa
%         bloccata da \( t_j \) solo se vale una certa condizione sui timestamp, ad esempio
%         \( TS_i < TS_j \), altrimenti \( t_i \) viene abortita e fatta ripartire con
%         lo stesso timestamp
%     \end{itemize}
%
%   \item \textbf{Rilevamento}: Si esegue un'analisi periodica della tabella di lock per
%     rilevare la presenza di un deadlock che corrisponde ad un ciclo nel grafo delle
%     relazioni di attesa. Questa tecnica è la più utilizzata nei sistemi reali.
% \end{itemize}
%
% \vspace{1em}
% \noindent
% Inoltre la scelta di gestire i lock in lettura come lock condivisi introduce il
% problema della \textbf{starvation} che tuttavia nel contesto delle DMBS risulta poco
% probabile, ma comunque risolvibile con le tecniche precedenti (in particolare il
% rilevamento).
%
% \begin{definition}[Starvation]
%   Se una risorsa \( x \) viene costantemente bloccata da una sequenza di transazioni 
%   che acquisiscono \( r\_lock \) su \( x \), un'eventuale transazione in attesa di
%   scrivere su \( x \) viene bloccata per un lungo periodo fino alla fine della
%   sequenza di letture.
%   \begin{figure}[H]
%     \centering
%     \includegraphics[width=0.9\textwidth]{starvation}
%     \caption{Esempio di starvation}
%   \end{figure}
% \end{definition}
%
% \subsubsection{Gestione della concorrenza in SQL}
% Siccome risulta molto pesante per il sistema gestire l'esecuzione concorrente di
% transazioni con una conseguente diminuzione delle prestazioni, SQL prevede la 
% possibilità di \textbf{rinunciare del tutto o in parte al controllo di concorrenza}
% per aumentare le prestazioni del sistema. È quindi possibile, a livello di una singola
% transazione, decidere di \textbf{tollerare alcune anomalie} di esecuzione concorrente.
% La seguente tabella mostra le possibili combinazioni di anomalie che si possono tollerare:
% \begin{table}[H]
%   \centering
%   \includegraphics[width=0.9\textwidth]{anomalie_tollerate}
%   \caption{Anomalie tollerate in SQL}
% \end{table}
% \textbf{Nota:}
% \begin{itemize}
%   \item Tutti i livelli richiedono il 2PL stretto per le scritture
%
%   \item \textbf{Serializable}: Richiede il 2PL stretto anche per le scritture e applica
%     il lock di predicato (cioè blocca tutte le risorse che soddisfano un certo
%     predicato) per evitare l'inserimento fantasma
%
%   \item \textbf{Repeatable read}: Applica il 2PL stretto per tutti i lock in lettura 
%     applicati a livello di tupla e non di tabella. Consente inserimenti e quindi non
%     evita l'anomalia di inserimento fantasma
%
%   \item \textbf{Read committed}: Richiede lock condivisi per le letture, ma non il
%     2PL stretto
%
%   \item \textbf{Read uncommitted}: Non applica lock in lettura
% \end{itemize}

\end{document}
